\section{Imagen gauge y composición hacia Yang--Mills}

\subsection{Estado gauge enriquecido y generadores de transferencia}

El objeto ya generado que recibe este capítulo es la estructura gauge completa
$\mathfrak G_{\rm gau}$; no constituye un dato externo ni una hipótesis terminal.
En cada nivel $m$ conserva la red
$\Lambda_m=a_m\mathbb Z^4$, $a_{m+1}=a_m/9$, la medida $\mu_m$, los
entrelazadores $J_m$, cuatro reflexiones $\Theta_{\nu,m}$, holonomía, color,
representación, acarreo de fusión, ruta, memoria, terminal y lector de curvatura.
Sus dos dinámicas se tipan con símbolos distintos antes de cualquier límite.

El semigrupo de una sola fibra de memoria es
\begin{equation}\label{eq:pdf2-ym-fibra-semigrupo}
 K^{\rm fib}_{m,t}
 =P_{\Omega_m}+e^{-3\kappa_{\rm av}t}(I-P_{\Omega_m})
 =e^{-tD_m^{\rm fib}},
 \qquad
 D_m^{\rm fib}:=3\kappa_{\rm av}(I-P_{\Omega_m}).
\end{equation}
Su espectro está contenido en $\{0,3\kappa_{\rm av}\}$. El generador de la
coordenatización producto completa es, en cambio,
\begin{equation}\label{eq:pdf2-ym-generador-producto}
 D_m^{\rm prod}:=3\kappa_{\rm av}
 \sum_{c\in\mathcal I_m}(I-E_{m,c}),
 \qquad
 K^{\rm prod}_{m,t}:=e^{-tD_m^{\rm prod}},
\end{equation}
donde las expectativas ortogonales $E_{m,c}$ conmutan. Si hay dos o más
coordenadas no constantes, $D_m^{\rm prod}$ posee niveles
$6\kappa_{\rm av},9\kappa_{\rm av},\ldots$; por ello
\eqref{eq:pdf2-ym-fibra-semigrupo} no puede declararse igual a
\eqref{eq:pdf2-ym-generador-producto}. El primero es la fibra simple; el segundo,
el producto de solapamiento.

\subsection{Teorema finito de Hoeffding, vacío y brecha alcanzada}

En la coordenatización producto construida en
\eqref{eq:pdf2-ym-owner-product-factorization}, los $E_{m,c}$ son expectativas
condicionales ortogonales y conmutantes, y la prueba de
\eqref{eq:pdf2-ym-owner-vacuum} establece
\begin{equation}\label{eq:pdf2-ym-interseccion-vacio}
 \bigcap_{c\in\mathcal I_m}\operatorname{Ran}E_{m,c}
 =\mathbb C\Omega_m.
\end{equation}
La descomposición de Hoeffding es la suma ortogonal
\[
 \mathcal H_m=\bigoplus_{S\subseteq\mathcal I_m}\mathcal H_{m,S},
 \qquad u=\sum_Su_S,
\]
donde $E_{m,c}u_S=0$ si $c\in S$ y $E_{m,c}u_S=u_S$ si $c\notin S$.
Entonces
\begin{equation}\label{eq:pdf2-ym-hoeffding}
 D_m^{\rm prod}u_S=3\kappa_{\rm av}|S|u_S,
 \qquad
 \langle u,D_m^{\rm prod}u\rangle
 =3\kappa_{\rm av}\sum_S|S|\|u_S\|^2.
\end{equation}
La componente $S=\varnothing$ es exactamente el vacío de
\eqref{eq:pdf2-ym-interseccion-vacio}. Por tanto
\begin{equation}\label{eq:pdf2-ym-cota-finita}
 D_m^{\rm prod}\succeq
 3\kappa_{\rm av}(I-P_{\Omega_m}),
 \qquad
 \|K^{\rm prod}_{m,t}(I-P_{\Omega_m})\|
 \le e^{-3\kappa_{\rm av}t}.
\end{equation}

En una coordenada de incidencia
$q=(q_1,\ldots,q_6)\in\mathbb F_3^6$, el testigo
\begin{equation}\label{eq:pdf2-ym-wc}
 W_c(q):=\mathbf1_{\{q_1+\cdots+q_6=0\}}-\frac13
\end{equation}
satisface
\begin{equation}\label{eq:pdf2-ym-wc-momentos}
 \mathbb EW_c=0,\qquad \|W_c\|^2=\frac29,\qquad
 D_m^{\rm prod}W_c=3\kappa_{\rm av}W_c.
\end{equation}
Así, en cada nivel finito, el umbral de
\eqref{eq:pdf2-ym-cota-finita} se alcanza. Este cálculo fija el espectro del
generador producto finito; la medida cuatridimensional continua de Yang--Mills se
construye materialmente en
\eqref{eq:pdf2-ym-gibbs-measure}--\eqref{eq:pdf2-ym-g-dynamics}.

\subsection{Compresión gauge física}

Sea $G$ compacto, conectado y simple. El grupo que actúa en el nivel $m$ es
\[
 \mathcal G_m^{\rm full}
 =G^{V(\Lambda_m)}\times
  \prod_{c\in\mathcal C_m}\mathbb F_3^4,
\]
donde el segundo factor actúa sobre las seis incidencias mediante una matriz
que debe quedar visible. Si las filas se indexan por
$(01,02,03,12,13,23)$ y las columnas por $0,1,2,3$, se fija
\begin{equation}\label{eq:pdf2-ym-incidence-matrix}
 B^+_{(ij),k}:=\delta_{ik}+\delta_{jk}\in\mathbb F_3,
 \qquad
 (B^+x)_{ij}=x_i+x_j.
\end{equation}
Esta es la incidencia no orientada del collar; no se confunde con el
coborde firmado de una arista. La acción finita es
\begin{equation}\label{eq:pdf2-ym-incidence-action}
 q_c\longmapsto q_c+B^+x_c,
 \qquad x_c\in\mathbb F_3^4.
\end{equation}
La proyección de Haar
\begin{equation}\label{eq:pdf2-ym-proyeccion-fisica}
 \mathbb P_m^{\rm phys}F
 :=\int_{\mathcal G_m^{\rm full}}F(g^{-1}\!\cdot\,)\,dg,
 \qquad
 \mathcal H_m^{\rm phys}:=\operatorname{Ran}\mathbb P_m^{\rm phys}
\end{equation}
es distinta de la marginal que olvida incidencia. La invariancia de las
expectativas $E_{m,c}$ da
\begin{equation}\label{eq:pdf2-ym-reduccion}
 [\mathbb P_m^{\rm phys},D_m^{\rm prod}]=0,
 \qquad
 [\mathbb P_m^{\rm phys},K_{m,t}^{\rm prod}]=0.
\end{equation}
Además, cada variable de $x_c\in\mathbb F_3^4$ aparece tres veces en la suma de
las seis incidencias, de modo que
\begin{equation}\label{eq:pdf2-ym-incidence-conservation}
 \sum_{i<j}(B^+x_c)_{ij}
 =3\sum_{i=0}^3x_{c,i}=0
 \quad\text{en }\mathbb F_3.
\end{equation}
Por ello $W_c$ es gauge
invariante y pertenece a $\mathcal H_m^{\rm phys}$. La brecha finita no es un
artefacto de una dilatación no física.

\subsection{Positividad por reflexión en cada nivel}

Para un kernel reversible común, defínase la medida de ocho caras
\begin{equation}\label{eq:pdf2-ym-ocho-caras}
 \Omega_m^{(8)}(dy_1\cdots dy_8)
 :=\int\mu_m(dz)\prod_{r=1}^{8}K^{\rm prod}_{m,a_m/2}(z,dy_r).
\end{equation}
Si $\Theta_{\nu,m}$ intercambia las dos mitades respecto de la dirección
$\nu$, la propiedad de Markov y la reversibilidad factorizan, para todo observable
cilíndrico $F$ soportado en el semiespacio positivo,
\begin{equation}\label{eq:pdf2-ym-os-finita}
 \int\overline{F(\Theta_{\nu,m}y)}F(y)\,d\Omega_m^{(8)}
 =\|\mathcal B_{m,\nu}F\|_2^2\ge0,
 \qquad \nu=1,\ldots,4.
\end{equation}
La marginal de caras opuestas recupera el mismo
$K^{\rm prod}_{m,a_m}$. De este modo las cuatro formas OS finitas son
consecuencia de una sola medida y un solo semigrupo, no de cuatro procesos
independientes.

\subsection{Estado gauge completo y coordenadas físicas de incidencia}

El propietario vigente no toma como espacio físico la marginal que conserva sólo
los enlaces. Su dominio es $\mathfrak G_{\rm gau}$, con enlace, marco,
representación, incidencia, color, ruta, orientación,
acarreo, memoria, supervivencia y terminal. La proyección a enlaces es un mapa de
olvido posterior. Juzgar el testigo físico después de aplicar ese olvido sustituye el
objeto y no transfiere la conclusión.

\subsubsection{Medida proyectiva y representación de enlaces}

Para una pantalla finita $\Gamma_m=(V_m,E_m)$ se parte del espacio levantado
\begin{equation}\label{eq:pdf2-ym-owner-lifted-space}
 \widetilde X_m:=G^{V_m}\times G^{E_m},\qquad
 d\widetilde\mu_m(h,z)
 =\prod_{v\in V_m}dh_v\prod_{e\in E_m}k_{\tau_{m,e}}(z_e)\,dz_e,
\end{equation}
y se expresa cada enlace mediante
\begin{equation}\label{eq:pdf2-ym-owner-link-publication}
 U_e=h_{s(e)}z_eh_{t(e)}^{-1}.
\end{equation}
Una arista gruesa $e$ se refina en la cadena ordenada
$e_1\cdots e_9$. Póngase
\(\boldsymbol\tau=(\tau_1,\ldots,\tau_9)\), con
\(\tau_j=\tau_{m+1,e_j}>0\) y
\(\tau=\tau_{m,e}=\sum_{j=1}^9\tau_j\). Sea
\(d\lambda_z^{(9)}\) la desintegración de Haar sobre la fibra del producto,
caracterizada por
\[
 \int_{G^9}F(z_1,\ldots,z_9)\,dz_1\cdots dz_9
 =\int_G\int_{z_1\cdots z_9=z}
   F\,d\lambda_z^{(9)}\,dz.
\]
El puente correcto para tiempos no necesariamente iguales es
\begin{equation}\label{eq:pdf2-ym-owner-heat-bridge}
 dB_{z,\boldsymbol\tau}^{(9)}(z_1,\ldots,z_9)
 =\frac{\prod_{j=1}^9k_{\tau_j}(z_j)}{k_\tau(z)}
   \,d\lambda_z^{(9)},
 \qquad z_1\cdots z_9=z,
\end{equation}
que es una probabilidad porque
\(k_{\tau_1}*\cdots*k_{\tau_9}=k_\tau\). En el refinamiento nonádico uniforme se
recupera \(\tau_j=\tau/9\), pero la fórmula no presupone esa igualdad.

Para hacer independiente la innovación sin borrar el producto grueso, se fija, para
cada $(e,z)$, un transporte boreliano
\begin{equation}\label{eq:pdf2-ym-owner-bridge-triangularization}
 \chi_{m,e,z}:\bigl((0,1)^{\mathbb N},du^{\otimes\mathbb N}\bigr)
 \longrightarrow
 \bigl(\{(z_j):z_1\cdots z_9=z\},B_{z,\boldsymbol\tau}^{(9)}\bigr)
\end{equation}
que es isomorfismo módulo conjuntos nulos. Existe porque ambas son realizaciones
estándar no atómicas; una elección queda fijada una vez y se transporta bajo
relabelados de aristas. La variable gruesa $z$ y la innovación $u^{\rm det}_{m+1,e}$
son entonces factores independientes en la coordenatización triangular. Los marcos nuevos se
integran con Haar. Ruta, acarreo, orientación y memoria se transportan por los
morfismos $\operatorname{Ext}_{\rm gau}$ del propio dato gauge.

La proyección gruesa multiplica los nueve incrementos y olvida exclusivamente las
innovaciones nuevas. Con esta definición se obtiene
\begin{equation}\label{eq:pdf2-ym-owner-projectivity}
 (\widehat\pi_{m+1,m})_*\widehat\mu_{m+1}^{\rm ov}
 =\widehat\mu_m^{\rm ov},\qquad
 \widehat J_mF=F\circ\widehat\pi_{m+1,m}.
\end{equation}
La misma arista aparece una sola vez en la medida; las seis caras incidentes son seis
lectores orientados de esa arista. Las rutas solapadas conservan el incremento
compartido y no se reemplazan por copias independientes.

\subsubsection{La fibra de incidencia es una coordenada física del estado completo}

La cola de incidencia incluida en $\mathfrak G_{\rm gau}$ se descompone sin perder
información:
\begin{equation}\label{eq:pdf2-ym-owner-seed-split}
 \Sigma:(\mathbb F_3^2)^{\mathbb N}\longrightarrow
 (\mathbb F_3^2)^{\mathbb N}\times\mathbb F_3^6,
 \qquad \Sigma_*\sigma=\sigma\otimes u_6.
\end{equation}
El primer factor continúa alimentando las coordenatizaciones de innovación y el segundo conserva
las seis incidencias $q_c=(q_{01},q_{02},q_{03},q_{12},q_{13},q_{23})$ del collar
$c$. En vez de dejar implícito qué coordenadas gobierna el generador, se fija una
coordenatización triangular recursiva
\begin{equation}\label{eq:pdf2-ym-owner-product-factorization}
 \widehat X_m
 =\widehat X_{m_0}^{\rm seed}
  \times\prod_{\ell=m_0+1}^{m}
       \mathcal U_{\ell/\ell-1}^{\rm det},
 \qquad
 \widehat\mu_m^{\rm ov}
 =\widehat\mu_{m_0}^{\rm seed}
  \otimes\bigotimes_{\ell=m_0+1}^{m}
       \upsilon_{\ell/\ell-1}^{\rm det}.
\end{equation}
El factor semilla enumera, una vez, los marcos y enlaces gruesos, las incidencias
$q_c$ y sus coordenatizaciones de marcación. Cada
$\mathcal U_{\ell/\ell-1}^{\rm det}$ enumera, una vez, los marcos nuevos, los puentes
\eqref{eq:pdf2-ym-owner-bridge-triangularization}, las incidencias nuevas y las
subcolas de marcación de ese refinamiento. La factorización es interna a
$\mathfrak G_{\rm gau}$ y no añade otro dato de partida.

Se registra el índice exhaustivo, disjunto y numerable
\begin{equation}\label{eq:pdf2-ym-owner-exhaustive-index}
 \mathcal I_m
 :=\mathcal I_{m_0}^{\rm frame}
 \sqcup\mathcal I_{m_0}^{\rm link}
 \sqcup\mathcal I_{m_0}^{q}
 \sqcup\mathcal I_{m_0}^{\rm mark}
 \sqcup\bigsqcup_{\ell=m_0+1}^{m}
 \bigl(
   \mathcal I_{\ell/\ell-1}^{\rm frame}
   \sqcup\mathcal I_{\ell/\ell-1}^{\rm bridge}
   \sqcup\mathcal I_{\ell/\ell-1}^{q}
   \sqcup\mathcal I_{\ell/\ell-1}^{\rm mark}
 \bigr).
\end{equation}
Cada factor independiente de \eqref{eq:pdf2-ym-owner-product-factorization} aparece
exactamente una vez; una subcola infinita se enumera por sus cilindros escalares. No
queda un factor genealógico fuera de $\mathcal I_m$.
En particular, las abreviaturas de la recurrencia significan literalmente
\begin{equation}\label{eq:pdf2-ym-owner-index-abbreviations}
 \begin{aligned}
 \mathcal I_{m_0}^{\rm seed}
 &:={\mathcal I}_{m_0}^{\rm frame}\sqcup
    {\mathcal I}_{m_0}^{\rm link}\sqcup
    {\mathcal I}_{m_0}^{q}\sqcup
    {\mathcal I}_{m_0}^{\rm mark},\\
 \mathcal I_{\ell/\ell-1}^{\rm det}
 &:={\mathcal I}_{\ell/\ell-1}^{\rm frame}\sqcup
    {\mathcal I}_{\ell/\ell-1}^{\rm bridge}\sqcup
    {\mathcal I}_{\ell/\ell-1}^{q}\sqcup
    {\mathcal I}_{\ell/\ell-1}^{\rm mark}.
 \end{aligned}
\end{equation}

La marginal de enlaces olvida $q_c$; la subálgebra física, en cambio, se obtiene
promediando la acción gauge completa
\begin{equation}\label{eq:pdf2-ym-full-gauge-group}
 \mathcal G_m^{\rm full}
 :=G^{V(\Lambda_m)}\times
 \prod_{c\in\mathcal C_m}\mathbb F_3^4,
 \qquad q_c\longmapsto q_c+B^+x_c.
\end{equation}
Por tanto la proyección física es
\begin{equation}\label{eq:pdf2-ym-full-physical-projection}
 \mathbb P_m^{\rm phys}F
 :=\int_{\mathcal G_m^{\rm full}}F(g^{-1}\!\cdot x)\,dg,
 \qquad
 \widehat{\mathscr H}_m^{\rm phys}
 :=\operatorname{Ran}\mathbb P_m^{\rm phys}.
\end{equation}
No coincide con el mapa de olvido
$\widehat X_m\to X_m^{\rm link}$. Esta distinción es el cortafuegos de tipo
que impide expulsar $W_c$ del sector físico.

\subsection{Generador gauge y circuito local de transporte}

La coordenatización triangular completa proporciona una unitaria
\begin{equation}\label{eq:pdf2-ym-owner-triangular-unitary}
 \widehat\Gamma_m:
 L^2(\widehat X_{m+1},\widehat\mu_{m+1}^{\rm ov})
 \longrightarrow
 L^2(\widehat X_m,\widehat\mu_m^{\rm ov})
 \widehat\otimes\mathcal K_{m+1}^{\rm det},
 \qquad
 \widehat\Gamma_m\widehat J_mF=F\otimes\mathbf1_{\rm det}.
\end{equation}
Para $\iota\in\mathcal I_m$, sea $E_{m,\iota}$ la esperanza que integra
exactamente el factor $\iota$ de
\eqref{eq:pdf2-ym-owner-product-factorization} y deja fijos todos los demás.
Son proyecciones Markov ortogonales, conservan la unidad y conmutan. Sobre el
núcleo denso $\mathscr C_m$ de cilindros que dependen de un número finito de
coordenadas se define la forma
\begin{equation}\label{eq:pdf2-ym-owner-full-form}
 \widehat{\mathcal E}_m[F]
 :=3\kappa_{\rm av}\sum_{\iota\in\mathcal I_m}
   \|(I-E_{m,\iota})F\|_2^2.
\end{equation}
La suma es finita para cada cilindro. Su clausura es una forma de Dirichlet y
determina el generador autoadjunto Markov
$\widehat H_m^{\rm ov}$. En un nivel semilla esto se escribe, en sentido de
formas,
\begin{equation}\label{eq:pdf2-ym-owner-base-generator}
 \widehat H_{m_0}^{\rm ov}
 =3\kappa_{\rm av}
  \sum_{\iota\in\mathcal I_{m_0}^{\rm seed}}(I-E_{m_0,\iota}),
\end{equation}
y cada refinamiento añade exactamente las expectativas de sus nuevas coordenadas:
\begin{equation}\label{eq:pdf2-ym-owner-refined-generator}
 \begin{aligned}
 \widehat H_{m+1}^{\rm ov}
 &=\widehat\Gamma_m^*
   (\widehat H_m^{\rm ov}\otimes I
    +I\otimes\widehat H_{m+1}^{\rm det})\widehat\Gamma_m,\\
 \widehat H_{m+1}^{\rm det}
 &=3\kappa_{\rm av}
   \sum_{\iota\in\mathcal I_{m+1/m}^{\rm det}}
   (I-E_{m+1,\iota}^{\rm det}).
 \end{aligned}
\end{equation}
La partición de colores no se deja nominal. Si una celda se indexa por
$n=(n_1,n_2,n_3,n_4)\in\mathbb Z^4$, se fija
\begin{equation}\label{eq:pdf2-ym-owner-729-coloring}
 \chi_m(n):=
 \bigl(n_1\bmod9,\ n_2\bmod9,\ n_3+3n_4\bmod9\bigr)
 \in(\mathbb Z/9\mathbb Z)^3,
 \qquad |(\mathbb Z/9\mathbb Z)^3|=9\cdot81.
\end{equation}
Si dos collares cerrados de radio celular uno se intersectan, su diferencia
$\delta$ satisface $|\delta_j|\le2$. La igualdad de colores fuerza
$\delta_1=\delta_2=0$ y
$\delta_3+3\delta_4\equiv0\pmod9$. Como el último entero pertenece a
$[-8,8]$, vale cero; las cotas $|\delta_3|,|\delta_4|\le2$ dan entonces
$\delta_3=\delta_4=0$. Por tanto, celdas distintas del mismo color tienen
collares disjuntos. Se define materialmente
\begin{equation}\label{eq:pdf2-ym-owner-729-blocks}
 \mathcal B_{m,r}:=\{\iota\in\mathcal I_m:
   \operatorname{cell}(\iota)=n, \chi_m(n)=r\},
 \qquad r\in(\mathbb Z/9\mathbb Z)^3.
\end{equation}
El transporte de una pantalla por $g\in E(4)$ transporta también esta
partición; no exige que una rotación preserve las coordenadas del rótulo.
Así, $\{\mathcal B_{m,r}\}_{r=1}^{9\cdot81}$ agrupa soportes espaciales
disjuntos;
el índice interno de cada bloque enumera sus factores de
\eqref{eq:pdf2-ym-owner-exhaustive-index}. Para un cilindro $F$, el producto sólo
contiene los bloques de los que depende $F$; la clausura define el producto fuerte
para toda la forma. El generador de un bloque conserva las multiplicidades de
Hoeffding; no se sustituye por un solo proyector que las borraría. Por
conmutatividad en cada coordenatización,
\begin{equation}\label{eq:pdf2-ym-owner-local-circuit}
\begin{aligned}
 H_{m,B}&:=3\kappa_{\rm av}\sum_{\iota\in B}(I-E_{m,\iota}),\\
 K_{m,B}&:=e^{-a_mH_{m,B}}
 =\prod_{\iota\in B}e^{-3\kappa_{\rm av}a_m(I-E_{m,\iota})},\\
 K_m^{\rm loc}&:=\mathop{\prod_{r,B}}^{\rm strong}K_{m,B}
 =e^{-a_m\widehat H_m^{\rm ov}}.
\end{aligned}
\end{equation}
El diámetro físico de cada bloque permanece uniformemente acotado. De aquí procede
la localidad: el generador no introduce una actualización instantánea sobre toda la
red.
La partición es compatible con Bianchi, no sólo con disjunción geométrica. Para una
celda cúbica $c$ se fija el producto orientado, con todos los factores transportados
al mismo vértice,
\begin{equation}\label{eq:pdf2-ym-owner-bianchi-block-product}
 \mathscr B_{m,c}(x):=
 \prod_{p\subset\partial c}^{\longrightarrow}
 T_p\operatorname{Hol}_p(x)T_p^{-1}.
\end{equation}
Cada arista interna de la frontera aparece dos veces y con orientaciones contrarias;
por tanto $\mathscr B_{m,c}=e$ punto a punto. Un bloque
$\mathcal B_{m,r}$ contiene el collar completo de toda celda que actualiza: si toca
una cara de $c$, actualiza simultáneamente el par inverso que comparte esa arista.
La cancelación anterior sigue siendo literal después de cada factor del circuito:
\begin{equation}\label{eq:pdf2-ym-owner-bianchi-block-invariance}
 K_{m,B}\mathbf1_{\{\mathscr B_{m,c}=e\}}
 =\mathbf1_{\{\mathscr B_{m,c}=e\}},
 \qquad
 K_m^{\rm loc}\mathbf1_{\{\mathscr B_{m,c}=e\}}
 =\mathbf1_{\{\mathscr B_{m,c}=e\}}.
\end{equation}
Así el calendario de $9\cdot81$ capas preserva Bianchi en cada paso y en el producto
fuerte, no sólo al final de una barrida.

La acción gauge permuta o traslada las coordenadas integradas por cada
$E_{m,\iota}$ y
preserva sus medidas. Por ello
\begin{equation}\label{eq:pdf2-ym-owner-physical-reduction}
 U_m(g)\widehat H_m^{\rm ov}=\widehat H_m^{\rm ov}U_m(g)
 \quad(g\in\mathcal G_m^{\rm full}),
 \qquad
 [\mathbb P_m^{\rm phys},\widehat H_m^{\rm ov}]=0,
 \qquad
 \mathbb P_{m+1}^{\rm phys}\widehat J_m
 =\widehat J_m\mathbb P_m^{\rm phys}.
\end{equation}
El semigrupo completo y su kernel Markov se denotan por
\begin{equation}\label{eq:pdf2-ym-owner-full-kernel}
 \widehat K_{m,t}^{\rm ov}:=e^{-t\widehat H_m^{\rm ov}},
 \qquad
\widehat K_{m,t}^{\rm ov}(x,dy)
 \quad(x,y\in\widehat X_m).
\end{equation}
La primera identidad de \eqref{eq:pdf2-ym-owner-physical-reduction} implica
\begin{equation}\label{eq:pdf2-ym-owner-full-kernel-equivariance}
 \widehat K_{m,t}^{\rm ov}(gx,gA)
 =\widehat K_{m,t}^{\rm ov}(x,A)
 \quad(g\in\mathcal G_m^{\rm full}).
\end{equation}
Para tipar la reducción sin usar un kernel restringido sobre el espacio equivocado,
sea $q_m:\widehat X_m\to Y_m:=\widehat X_m/\mathcal G_m^{\rm full}$ el mapa de
órbitas y $\nu_m=(q_m)_*\widehat\mu_m^{\rm ov}$. El pullback
\begin{equation}\label{eq:pdf2-ym-owner-quotient-unitary}
 \mathcal W_m:L^2(Y_m,\nu_m)\longrightarrow
 \widehat{\mathscr H}_m^{\rm phys},
 \qquad \mathcal W_m\phi=\phi\circ q_m
\end{equation}
es unitario. La equivariancia del kernel completo permite definir
\begin{equation}\label{eq:pdf2-ym-owner-physical-kernel}
 K_{m,t}^{\rm phys}(q_mx,A)
 :=\widehat K_{m,t}^{\rm ov}(x,q_m^{-1}A),
 \qquad A\subseteq Y_m,
\end{equation}
independientemente del representante $x$. El objeto físico queda entonces tipado por
\begin{equation}\label{eq:pdf2-ym-owner-D}
 D_m^{\rm YM}
 :=\left.\widehat H_m^{\rm ov}\right|_{\widehat{\mathscr H}_m^{\rm phys}},
 \qquad
 e^{-tD_m^{\rm YM}}
 =\mathcal W_mK_{m,t}^{\rm phys}\mathcal W_m^{-1}.
\end{equation}
Éste es el generador producto completo; no es el generador de una sola fibra de
\eqref{eq:pdf2-ym-fibra-semigrupo}.

\subsection{Acción de Yang--Mills, medida dependiente de $g$ y transporte exacto}

Hasta aquí se ha escrito la coordenatización cinemática de referencia. La denotaremos
$\widehat\mu_m^0:=\widehat\mu_m^{\rm ov}$ y
$\widehat H_m^0:=\widehat H_m^{\rm ov}$, con
$\widehat K_{m,t}^0:=e^{-t\widehat H_m^0}=\widehat K_{m,t}^{\rm ov}$.
Para convertir el lector estructural
$\mathcal P_m^{F^2}$ en una ley, y no en una comparación posterior, se descarga
primero su raíz positiva sobre la factorización triangular. La polarización de
$\mathcal P_m^{F^2}$ en un collar $b$ es el Gram positivo
\[
 Q_{m,b}(u,v):=\frac14\bigl(
 \mathcal P_{m,b}^{F^2}(u+v)-\mathcal P_{m,b}^{F^2}(u-v)
 \bigr).
\]
Se divide el espacio finito de lectores del collar por $\ker Q_{m,b}$, se completa
con $Q_{m,b}$ y se denota por $\mathcal E_{m,b}^{F}$ el Hilbert resultante. Si
$\zeta_{m,b}(x_b)$ es el vector de lectores centrados de la coordenada triangular
$x_b$, se fija materialmente
\begin{equation}\label{eq:pdf2-ym-action-root}
 \xi_{m,b}(x_b):=Q_{m,b}^{1/2}\zeta_{m,b}(x_b)
 \in\mathcal E_{m,b}^{F},
 \qquad
 p_{m,b}(x_b):=\|\xi_{m,b}(x_b)\|^2.
\end{equation}
El índice $b$ recorre los collares Gram de soporte mínimo. Un collar se asigna
al primer color de \eqref{eq:pdf2-ym-owner-729-coloring} que contiene su celda base;
así cada término aparece una vez y los collares de un mismo color son disjuntos.
Más precisamente, la diagonalización positiva agrupa los factores triangulares en
una partición disjunta
\begin{equation}\label{eq:pdf2-ym-action-factor-partition}
 \mathcal I_m^{F}
 =\bigsqcup_{r\in(\mathbb Z/9\mathbb Z)^3}
   \bigsqcup_{b\in\mathcal B_{m,r}^{F}}\mathcal I_{m,b},
 \qquad x_b:=(x_\iota)_{\iota\in\mathcal I_{m,b}},
 \qquad
 \mathcal B_{m,r}^{F}:=\{b:Q_{m,b}\ne0,\ \chi_m(b)=r\}.
\end{equation}
Ninguna coordenada triangular pertenece a dos $\mathcal I_{m,b}$; si un lector
geométrico toca dos collares, su modo de Gram se asigna al primero y su complemento
ortogonal al segundo. Ésta es la descomposición espectral finita de la forma positiva,
no una duplicación de la variable física.
La raíz de la suma ortogonal, calculada color por color, da la identidad, no una
aproximación,
\begin{equation}\label{eq:pdf2-ym-action-reader-sum}
 \mathcal P_m^{F^2}(x)
 =\sum_{r\in(\mathbb Z/9\mathbb Z)^3}
   \sum_{b\in\mathcal B_{m,r}^{F}}p_{m,b}(x_b),
 \qquad
 S_{m,g}(x):=\frac1{4g^2}\mathcal P_m^{F^2}(x),
 \quad g\in(0,\infty).
\end{equation}
Aquí $x_b:\widehat X_m\to X_{m,b}$ es la proyección de coordenadas del
collar, de modo que dominio, acción y codominio de cada sumando quedan fijados.
El lector es gauge invariante porque $Q_{m,b}$ usa el producto invariante de
$\mathfrak g$; es invariante por reflexión porque una reflexión sólo permuta los
seis pares orientados; y es natural por relabelado euclídeo.
La partición \eqref{eq:pdf2-ym-action-factor-partition} y la coordenatización producto dan
explícitamente
\begin{equation}\label{eq:pdf2-ym-action-factor-measures}
 \widehat\mu_m^0=\lambda_m^{\rm inert}\otimes
   \bigotimes_{r,b}\lambda_{m,b}^0,
 \qquad
 d\lambda_{m,b}^g
 :=(Z_{m,b}^g)^{-1}e^{-p_{m,b}/(4g^2)}d\lambda_{m,b}^0,
 \qquad
 \widehat\mu_{m,g}^{\rm YM}=\lambda_m^{\rm inert}\otimes
   \bigotimes_{r,b}\lambda_{m,b}^g.
\end{equation}
El factor inerte contiene exactamente las coordenadas sobre las que el lector es
cero; no se elimina de la teoría.

La separación grueso--detalle de
\eqref{eq:pdf2-ym-owner-product-factorization} es ortogonal para la misma raíz:
la polarización directa da
$Q_{m+1}=Q_m\oplus Q_{m+1/m}^{\rm det}$ y
$\zeta_{m+1}=\zeta_m\oplus\zeta_{m+1/m}^{\rm det}$.
Si $x_{m+1}=\widehat\Gamma_m^{-1}(x_m,u)$, entonces
\begin{equation}\label{eq:pdf2-ym-action-cocycle}
 \xi_{m+1}(x_{m+1})
 =\xi_m(x_m)\oplus\xi_{m+1/m}^{\rm det}(u),
 \qquad
 S_{m+1,g}\circ\widehat\Gamma_m^{-1}
 =S_{m,g}+S_{m+1/m,g}^{\rm det}.
\end{equation}
Esta es la forma precisa en que carry y memoria impiden contar dos veces una cara
ancestral: la parte gruesa no se vuelve a sumar en el detalle. En particular,
$Z_{m+1,g}=Z_{m,g}Z_{m+1/m,g}^{\rm det}$ para los normalizadores y se obtiene la
colección de probabilidades genuinamente dependiente del acoplamiento
\begin{equation}\label{eq:pdf2-ym-gibbs-measure}
 d\widehat\mu_{m,g}^{\rm YM}(x)
 :=Z_{m,g}^{-1}e^{-S_{m,g}(x)}d\widehat\mu_m^0(x),
 \qquad
 (\widehat\pi_{m+1,m})_*\widehat\mu_{m+1,g}^{\rm YM}
 =\widehat\mu_{m,g}^{\rm YM}.
\end{equation}
Para escribir correctamente Schwinger--Dyson respecto de Haar, y no fingir que la
medida de referencia es plana, sea $\lambda_m$ el producto de Haar, Lebesgue y
conteo uniforme de la coordenatización levantada, y sea
$R_m:=d\widehat\mu_m^0/d\lambda_m>0$. La acción total regulada es
\begin{equation}\label{eq:pdf2-ym-total-regulated-action}
 \mathscr S_{m,g}:=S_{m,g}-\log R_m,
 \qquad
 d\widehat\mu_{m,g}^{\rm YM}
 =\mathscr Z_{m,g}^{-1}e^{-\mathscr S_{m,g}}d\lambda_m.
\end{equation}
El término $-\log R_m$ contiene los kernels de calor y las coordenatizaciones de puente de la
estructura gauge de partida; es independiente de $g$. El término de Yang--Mills que
cambia el acoplamiento y cuyo límite clásico se calcula abajo es exactamente
$S_{m,g}=\mathcal P_m^{F^2}/(4g^2)$.
La consistencia de la segunda igualdad y la estandaridad de los espacios finitos
permiten aplicar Kolmogórov y definen la única medida cilíndrica
$\widehat\mu_{\infty,g}^{\rm YM}$ cuyas marginales son
$\widehat\mu_{m,g}^{\rm YM}$.
No se ha usado aquí ningún autovalor: primero se han fijado $g$, la acción y
la medida.

Para definir la dinámica sobre esa medida sin perder producto, gauge ni las
identidades superficiales, se construye el transporte, no se lo postula. En cada
factor no atómico $X_{m,b}$ se fija una coordenatización boreliana
$\rho_{m,b}:X_{m,b}/\mathcal G_{m,b}\to(0,1)$ y se desintegra primero sobre el
cociente gauge y después sobre la órbita de Haar. Sean $F_{m,b}^0$ y
$F_{m,b}^g$ las funciones de distribución continuas de $\rho_{m,b}$ para la
medida de referencia y la medida reponderada. El transporte triangular es
\begin{equation}\label{eq:pdf2-ym-gibbs-transport-factor}
 t_{m,b,g}:=(F_{m,b}^g)^{-1}\circ F_{m,b}^0,
 \qquad
 T_{m,b,g}(y,h):=(t_{m,b,g}(y),h),
 \qquad
 (T_{m,b,g})_*\lambda_{m,b}^0=\lambda_{m,b}^g.
\end{equation}
En un factor puramente atómico sobre el que $p_{m,b}=0$ se toma
$T_{m,b,g}=I$. En presencia de átomos y una coordenada continua, la coordenatización es el
orden lexicográfico átomo--intervalo y la misma fórmula se aplica al espacio
estándar no atómico total. Los transportes de los $9\cdot81$ colores se componen
en el orden fijado por \eqref{eq:pdf2-ym-owner-729-coloring}; dentro de un color
conmutan porque los collares son disjuntos. Se obtiene
\begin{equation}\label{eq:pdf2-ym-gibbs-transport-global}
 T_{m,g}:=\prod_{r\in(\mathbb Z/9\mathbb Z)^3}
           \prod_{b\in\mathcal B_{m,r}^{F}}T_{m,b,g},
 \quad
 (T_{m,g})_*\widehat\mu_m^0=\widehat\mu_{m,g}^{\rm YM},
 \quad
 \widehat\pi_{m+1,m}T_{m+1,g}=T_{m,g}\widehat\pi_{m+1,m}.
\end{equation}
La última igualdad se comprueba factor a factor usando
\eqref{eq:pdf2-ym-action-cocycle}: el transporte del factor grueso es el ya fijado y
los transportes nuevos actúan sólo en la innovación. Las coordenatizaciones de cociente se
eligen en una pantalla representante y se trasladan; por ello $T_{m,g}$ conmuta con
gauge y reflexión y, para $a\in E(4)$, satisface
$T_{a\alpha,g}=U_\alpha(a)T_{\alpha,g}U_\alpha(a)^{-1}$.

La composición
\begin{equation}\label{eq:pdf2-ym-full-probability-isomorphism}
 \mathcal T_{m,g}:L^\infty(\widehat X_m,\widehat\mu_{m,g}^{\rm YM})
 \longrightarrow L^\infty(\widehat X_m,\widehat\mu_m^0),
 \qquad \mathcal T_{m,g}F:=F\circ T_{m,g}
\end{equation}
es, por construcción, un isomorfismo de espacios de probabilidad y de
$*$-álgebras, y se prolonga unitariamente a $L^2$. En particular,
\begin{equation}\label{eq:pdf2-ym-full-algebra-identities}
 \mathcal T_{m,g}(FG)=\mathcal T_{m,g}F\,\mathcal T_{m,g}G,
 \quad
 \mathcal T_{m,g}(F^*)=(\mathcal T_{m,g}F)^*,
 \quad
 \int F\,d\widehat\mu_{m,g}^{\rm YM}
 =\int\mathcal T_{m,g}F\,d\widehat\mu_m^0.
\end{equation}
El dominio de cada $T_{m,b,g}$ es el propio espacio de holonomías que satisface
la ley de subdivisión y la restricción de Bianchi. Por tanto su imagen permanece
en ese espacio. Aplicando la multiplicatividad anterior generador a generador se
obtienen las identidades puntuales
\begin{equation}\label{eq:pdf2-ym-transport-surface-bianchi}
 \begin{aligned}
 \mathcal T_{m,g}(\operatorname{Hol}_p)
 &=\prod_{p'\subset p}^{\longrightarrow}
   \mathcal T_{m,g}(T_{p'}\operatorname{Hol}_{p'}T_{p'}^{-1}),\\
 \mathcal T_{m,g}(F_{{\rm cl},m})
 &=F_{{\rm cl},m}\circ T_{m,g},\\
 \prod_{p\subset\partial c}^{\longrightarrow}
 \mathcal T_{m,g}(T_p\operatorname{Hol}_pT_p^{-1})&=e,
 \qquad
 \operatorname{Alt}_{\lambda\mu\nu}
 D_\lambda(F_{{\rm cl},m}\circ T_{m,g})_{\mu\nu}=0.
 \end{aligned}
\end{equation}
La tercera línea es la cancelación arista por arista de la frontera orientada de
la frontera de $c$; la cuarta es su polarización infinitesimal. Así el transporte
preserva productos, estado, clover y Bianchi, no sólo el primer caos.

Finalmente se transportan expectativas, generador, kernel, proyección e inclusión:
\begin{equation}\label{eq:pdf2-ym-g-dynamics}
 \begin{aligned}
 E_{m,\iota}^{g}&:=\mathcal T_{m,g}^{-1}E_{m,\iota}\mathcal T_{m,g},&
 \widehat H_{m,g}^{\rm YM}&:=\mathcal T_{m,g}^{-1}\widehat H_m^0\mathcal T_{m,g},\\
 \widehat K_{m,t}^{g}&:=\mathcal T_{m,g}^{-1}\widehat K_{m,t}^{0}\mathcal T_{m,g},&
 \mathbb P_{m,g}^{\rm phys}&:=\mathcal T_{m,g}^{-1}
     \mathbb P_m^{\rm phys}\mathcal T_{m,g},\\
 \widehat J_{m,g}&:=\mathcal T_{m+1,g}^{-1}
     \widehat J_m\mathcal T_{m,g},&
 D_{m,g}^{\rm YM}&:=\left.\widehat H_{m,g}^{\rm YM}
     \right|_{\operatorname{Ran}\mathbb P_{m,g}^{\rm phys}}.
 \end{aligned}
\end{equation}
La primera línea no es sólo una conjugación abstracta. Usando
\eqref{eq:pdf2-ym-action-factor-measures}, su acción sobre un cilindro es la
actualización de baño térmico de la acción:
\begin{equation}\label{eq:pdf2-ym-gibbs-conditional-expectation}
 (E_{m,b}^{g}F)(x_{\ne b})
 =\frac{\displaystyle
   \int_{X_{m,b}}F(x_{\ne b},u)e^{-p_{m,b}(u)/(4g^2)}
       d\lambda_{m,b}^0(u)}{\displaystyle
   \int_{X_{m,b}}e^{-p_{m,b}(u)/(4g^2)}d\lambda_{m,b}^0(u)},
 \qquad
 \widehat H_{m,g}^{\rm YM}
 =3\kappa_{\rm av}\sum_b(I-E_{m,b}^{g}).
\end{equation}
Así tanto la medida como cada transición local contienen $g$ y la acción
$p_{m,b}/(4g^2)$ de forma explícita.
Son Markov porque $\mathcal T_{m,g}\mathbf1=\mathbf1$, locales porque cada
$T_{m,b,g}$ tiene el mismo collar que $p_{m,b}$, y dependen de $g$ por
\eqref{eq:pdf2-ym-gibbs-transport-factor}. La equivalencia es espectral y no
selecciona la medida a partir del espectro:
Para cerrar además las ecuaciones dinámicas, sea $X_{m,b}^a$ el campo
izquierdo-invariante en el factor de enlace $b$ asociado a
$e_a\in\mathfrak g$; en una coordenatización continua de innovación se usa la derivada
direccional ordinaria, y las coordenadas finitas se tratan por diferencia exacta.
Para todo cilindro suave $F$ el dominio común es
$\mathscr C_m^1:=\bigcap_{b,a}\operatorname{Dom}X_{m,b}^a$, tomando soporte
compacto en las coordenatizaciones abiertas de innovación; en los factores de grupo no hay
frontera. La invariancia de Haar y la regla del producto calculan
\begin{equation}\label{eq:pdf2-ym-schwinger-dyson-finite}
 \int X_{m,b}^aF\,d\widehat\mu_{m,g}^{\rm YM}
 =\int F\,X_{m,b}^a\mathscr S_{m,g}\,d\widehat\mu_{m,g}^{\rm YM},
 \qquad F\in\mathscr C_m^1.
\end{equation}
En una coordenada finita, para
$\Delta_\varepsilon F(q):=F(q+\varepsilon)-F(q)$, el cambio de variable
$q\mapsto q+\varepsilon$ da la versión exacta, sin derivada formal,
\begin{equation}\label{eq:pdf2-ym-schwinger-dyson-discrete}
 \int\Delta_\varepsilon F(q)\,d\widehat\mu_{m,g}^{\rm YM}
 =\int F(q)
  \left(e^{\mathscr S_{m,g}(q)-\mathscr S_{m,g}(q-\varepsilon)}-1\right)
  d\widehat\mu_{m,g}^{\rm YM}(q).
\end{equation}
Si $\delta_\varepsilon$ es la suma orientada de esos campos en las aristas que
inciden en un vértice, la invariancia gauge de la medida de referencia y de
$S_{m,g}$ da $\delta_\varepsilon\mathscr S_{m,g}=0$ y, por tanto, la identidad de Ward
\begin{equation}\label{eq:pdf2-ym-ward-finite}
 \int\delta_\varepsilon F\,d\widehat\mu_{m,g}^{\rm YM}=0.
\end{equation}
Para el factor gauge finito $\mathbb F_3^4$, la invariancia
$\mathscr S_{m,g}(q+\varepsilon)=\mathscr S_{m,g}(q)$ reduce
\eqref{eq:pdf2-ym-schwinger-dyson-discrete} a la misma identidad Ward con
$\delta_\varepsilon$ reemplazada por $\Delta_\varepsilon$.
Las dos expresiones son compatibles con $\widehat J_{m,g}$ porque
\eqref{eq:pdf2-ym-action-cocycle} separa las derivadas gruesas de las de detalle.
Todo cilindro del límite reside en algún nivel; así
\eqref{eq:pdf2-ym-schwinger-dyson-finite} y
\eqref{eq:pdf2-ym-ward-finite} pasan literalmente al estado cofinal. Ni la
ecuación Schwinger--Dyson ni Ward contienen la brecha como premisa.

\begin{equation}\label{eq:pdf2-ym-g-gap-before-limit}
 D_{m,g}^{\rm YM}\succeq3\kappa_{\rm av}
 (I-P_{\Omega_{m,g}}),
 \qquad
 W_{c,g}:=\mathcal T_{m,g}^{-1}W_c,
 \qquad
 D_{m,g}^{\rm YM}W_{c,g}=3\kappa_{\rm av}W_{c,g}.
\end{equation}
En lo que sigue se fija un $g\in(0,\infty)$ y se suprime ese subíndice; las
medidas, expectativas, inclusiones y generadores sin superíndice son los objetos
de \eqref{eq:pdf2-ym-g-dynamics}. Las holonomías geométricas permanecen como
funciones coordenadas canónicas sobre el espacio de configuraciones y, por ello, su
ley sí cambia con $g$; $\mathcal T_{m,g}$ las lleva a funciones compuestas de la
coordenatización de referencia. Sólo el testigo espectral se transporta explícitamente en
\eqref{eq:pdf2-ym-g-gap-before-limit}. Esta convención no borra $g$: su residencia explícita es
\eqref{eq:pdf2-ym-gibbs-measure}--\eqref{eq:pdf2-ym-g-dynamics}.

\subsection{Vacío simple y testigo que alcanza la brecha}

La intersección de todas las expectativas del estado completo es la constante:
\begin{equation}\label{eq:pdf2-ym-owner-vacuum}
 \bigcap_{\iota\in\mathcal I_m}\operatorname{Ran}E_{m,\iota}
 =\mathbb C\Omega_m.
\end{equation}
Esta identidad no se introduce como hipótesis. En la coordenatización producto
\eqref{eq:pdf2-ym-owner-product-factorization}, cada $E_{m,\iota}$ integra una
coordenada del índice exhaustivo
\eqref{eq:pdf2-ym-owner-exhaustive-index} y deja fijas las restantes. Si
$u$ pertenece a la intersección, sus expectativas condicionales respecto de todo
cilindro finito son constantes. La convergencia martingala de esos cilindros a la
$\sigma$-álgebra completa da que $u$ es constante casi en todas partes. La inclusión
inversa es inmediata porque toda expectativa conserva la unidad. Esto prueba
\eqref{eq:pdf2-ym-owner-vacuum} y, después de la proyección gauge, deja un vacío
físico unidimensional.
La descomposición de Hoeffding de
\eqref{eq:pdf2-ym-owner-base-generator}--
\eqref{eq:pdf2-ym-owner-refined-generator} da
\begin{equation}\label{eq:pdf2-ym-owner-hoeffding}
 \widehat{\mathcal E}_m(u)
 =3\kappa_{\rm av}\sum_{S\Subset\mathcal I_m}|S|\,\|u_S\|^2,
 \qquad
 D_m^{\rm YM}\succeq
 3\kappa_{\rm av}(I-P_{\Omega_m}).
\end{equation}
La función
\begin{equation}\label{eq:pdf2-ym-owner-wc}
 W_c(q):=\mathbf1_{\{q_1+\cdots+q_6=0\}}-\frac13
\end{equation}
es invariante bajo $G^{V(\Lambda_m)}$ porque no depende de enlaces ni marcos. Para
la acción de incidencia, cada $x_i$ aparece tres veces, luego
\begin{equation}\label{eq:pdf2-ym-owner-wc-invariance}
 \sum_{i<j}(B^+x)_{ij}=3\sum_i x_i=0
 \quad\text{en }\mathbb F_3.
\end{equation}
Así $\mathbb P_m^{\rm phys}W_c=W_c$. Bajo $u_6$, la suma de los seis trits es
uniforme; por cálculo directo,
\begin{equation}\label{eq:pdf2-ym-owner-wc-moments}
 \mathbb EW_c=0,\qquad
 \|W_c\|^2=\frac29,\qquad
 \mathbb E(W_c^3)=\frac2{27}.
\end{equation}
Sólo el proyector $I-E_{m,q_c}$ actúa no trivialmente sobre $W_c$. Por tanto
\begin{equation}\label{eq:pdf2-ym-owner-wc-eigenvalue}
 D_m^{\rm YM}W_c=3\kappa_{\rm av}W_c.
\end{equation}
La cota inferior de \eqref{eq:pdf2-ym-owner-hoeffding} no es sólo una cota: el
complemento físico contiene un vector no nulo que alcanza el umbral.

\subsection{Refinamiento nonádico y persistencia en el límite}

Las inclusiones físicas satisfacen
\begin{equation}\label{eq:pdf2-ym-owner-intertwining}
 D_{m+1}^{\rm YM}\widehat J_m
 =\widehat J_mD_m^{\rm YM},\qquad
 \bigl(e^{-a_{m+1}D_{m+1}^{\rm YM}}\bigr)^9\widehat J_m
 =\widehat J_me^{-a_mD_m^{\rm YM}},\qquad a_{m+1}=a_m/9.
\end{equation}
En el límite inductivo
\begin{equation}\label{eq:pdf2-ym-owner-inductive-limit}
 \widehat{\mathscr H}_\infty^{\rm phys}
 =\varinjlim(\widehat{\mathscr H}_m^{\rm phys},\widehat J_m)
\end{equation}
se define sobre la unión algebraica
\begin{equation}\label{eq:pdf2-ym-owner-limit-form}
 \mathcal E_\infty[\widehat J_{m,\infty}u]
 :=\langle u,D_m^{\rm YM}u\rangle.
\end{equation}
Su generador se denota
$D_{\infty,g}^{\rm YM}$; de acuerdo con la convención posterior a
\eqref{eq:pdf2-ym-g-gap-before-limit}, se escribe también
$D_\infty^{\rm YM}:=D_{\infty,g}^{\rm YM}$.
El primer entrelazamiento de \eqref{eq:pdf2-ym-owner-intertwining} hace que la
definición sea independiente del representante. Las truncaciones de Hoeffding son
un núcleo de forma y dan una sucesión de recuperación; Fatou en la suma no negativa
da la semicontinuidad inferior. La clausura posee, pues,
\begin{equation}\label{eq:pdf2-ym-owner-limit-gap}
 \overline{\mathcal E}_\infty[u]
 \ge3\kappa_{\rm av}\|(I-P_{\Omega_\infty})u\|^2.
\end{equation}
La inclusión conserva las coordenadas ancestrales, de modo que
$\widehat J_{m,\infty}W_c$ sigue siendo no nulo y satisface
\begin{equation}\label{eq:pdf2-ym-owner-limit-witness}
 D_\infty^{\rm YM}\widehat J_{m,\infty}W_c
 =3\kappa_{\rm av}\widehat J_{m,\infty}W_c.
\end{equation}
En consecuencia,
\begin{equation}\label{eq:pdf2-ym-owner-exact-gap}
 \inf\bigl(\operatorname{Spec}D_\infty^{\rm YM}\setminus\{0\}\bigr)
 =3\kappa_{\rm av}>0,
 \qquad \ker D_\infty^{\rm YM}=\mathbb C\Omega_\infty.
\end{equation}

\subsection{Cuatro reconstrucciones OS de la misma dinámica}

El kernel completo de \eqref{eq:pdf2-ym-owner-full-kernel} define una ley reversible
sobre $\widehat X_m$. Para la reconstrucción física se usa su descenso bien tipado
\eqref{eq:pdf2-ym-owner-physical-kernel}: sobre el cociente $Y_m$ se define
\begin{equation}\label{eq:pdf2-ym-owner-eight-face-law}
 \Omega_{m,{\rm phys}}^{(8)}(d\bar y_1\cdots d\bar y_8)
 =\int\nu_m(d\bar z)
  \prod_{r=1}^8K_{m,a_m/2}^{\rm phys}(\bar z,d\bar y_r).
\end{equation}
Para cada dirección $\nu=1,\ldots,4$, la reflexión intercambia las cuatro caras
positivas con las negativas. Condicionando respecto del estado central $z$ y usando
reversibilidad,
\begin{equation}\label{eq:pdf2-ym-owner-four-os}
 \int\overline{F(\Theta_{\nu,m}y)}F(y)\,
 d\Omega_{m,{\rm phys}}^{(8)}(y)
 =\|\mathcal B_{m,\nu}F\|_2^2\ge0.
\end{equation}
La marginal de la pareja opuesta normal a $\nu$ es
\begin{equation}\label{eq:pdf2-ym-owner-os-marginal}
 \nu_m(d\bar x)K_{m,a_m}^{\rm phys}(\bar x,d\bar y).
\end{equation}
Para identificar el Hamiltoniano de una dirección se toma el subespacio OS de
observables de una cara positiva simple,
$F(y)=F_\nu(\bar y_{\nu,+})$; los observables de varias caras conservan la
positividad anterior, pero no se confunden con este espacio de transferencia. En el
cociente de cara simple, el mapa
\begin{equation}\label{eq:pdf2-ym-owner-os-unitary}
 \mathcal V_{m,\nu}[F]
 :=e^{-a_mD_m^{\rm YM}/2}\mathcal W_mF
\end{equation}
es isométrico para la norma OS y tiene rango denso. Su extensión identifica el
cociente OS con $\widehat{\mathscr H}_m^{\rm phys}$. Póngase
$\overline D_m:=\mathcal W_m^{-1}D_m^{\rm YM}\mathcal W_m$ y
$\overline J_m:=\mathcal W_{m+1}^{-1}\widehat J_m\mathcal W_m$. El
entrelazamiento y $a_m=9a_{m+1}$ dan
\begin{equation}\label{eq:pdf2-ym-owner-os-range-inclusion}
 \overline J_me^{-a_m\overline D_m/2}
 =e^{-a_m\overline D_{m+1}/2}\overline J_m
 =\bigl(e^{-a_{m+1}\overline D_{m+1}/2}\bigr)^9\overline J_m.
\end{equation}
El último miembro pertenece al rango de
$e^{-a_{m+1}\overline D_{m+1}/2}$; por eso el inverso de $\mathcal V_{m+1,\nu}$
que aparece a continuación se aplica sólo en su rango y está bien definido. Los
conectores
\begin{equation}\label{eq:pdf2-ym-owner-os-connectors}
 \mathcal J_{m,\nu}^{\rm OS}
 :=\mathcal V_{m+1,\nu}^{-1}\widehat J_m\mathcal V_{m,\nu}
\end{equation}
son isometrías y entrelazan los cuatro Hamiltonianos reconstruidos:
\begin{equation}\label{eq:pdf2-ym-owner-common-os-hamiltonian}
 H_{{\rm OS},m+1}^{(\nu)}\mathcal J_{m,\nu}^{\rm OS}
 =\mathcal J_{m,\nu}^{\rm OS}H_{{\rm OS},m}^{(\nu)},
 \qquad
 \mathcal V_{\infty,\nu}H_{{\rm OS},\infty}^{(\nu)}
 \mathcal V_{\infty,\nu}^{-1}=D_\infty^{\rm YM}.
\end{equation}
Las cuatro reflexiones no fabrican cuatro medidas: son cuatro representaciones del mismo
estado, del mismo kernel y del mismo generador.

\subsection{Localidad, acción euclídea y representación de curvatura}

La categoría dirigida de pantallas incluye refinamientos, overlays y transportes
euclídeos. Si $a\in E(4)$ y $\alpha$ es una pantalla, el relabelado completo define
\begin{equation}\label{eq:pdf2-ym-owner-e4-finite-transport}
 U_\alpha(a):\widehat{\mathscr H}_\alpha^{\rm phys}
 \longrightarrow\widehat{\mathscr H}_{a\alpha}^{\rm phys},
 \qquad
 U_\beta(a)\widehat J_{\alpha\beta}
 =\widehat J_{a\alpha,a\beta}U_\alpha(a).
\end{equation}
El transporte preserva medida, generador, ideales nulos y soportes. La identidad de
naturalidad permite definir, sin elegir representante, una unitaria sobre el colímite
algebraico:
\begin{equation}\label{eq:pdf2-ym-owner-e4-colimit-action}
 U(a)\widehat J_{\alpha,\infty}F
 :=\widehat J_{a\alpha,\infty}U_\alpha(a)F.
\end{equation}
Los relabelados satisfacen $U_{a\alpha}(b)U_\alpha(a)=U_\alpha(ba)$ y preservan
norma; por ello \eqref{eq:pdf2-ym-owner-e4-colimit-action} se prolonga a una
representación unitaria de $E(4)$.

\subsubsection{Curvatura local del canal de incidencia}

Todo grupo compacto, conectado y simple posee un toro maximal no trivial. Se fija
$t_3\in G$ de orden exactamente tres. Para un collar $c$, sea
\begin{equation}\label{eq:pdf2-ym-owner-incidence-flux}
 \omega_c(q):=\sum_{i<j}q_{ij}\in\mathbb F_3.
\end{equation}
La identidad \eqref{eq:pdf2-ym-incidence-conservation} prueba que $\omega_c$ es
invariante bajo la acción finita. Si $h_c$ es el marco en el vértice base del collar,
se define la holonomía geométrica y, para el testigo espectral, su imagen transportada
en la ley a acoplamiento $g$:
\begin{equation}\label{eq:pdf2-ym-owner-incidence-holonomy}
 \operatorname{Hol}_{m,c}^{\rm inc,geom}:=h_ct_3^{\,\omega_c(q)}h_c^{-1}\in G,
 \qquad
 \operatorname{Hol}_{m,c,g}^{\rm inc,spec}
 :=\mathcal T_{m,g}^{-1}\operatorname{Hol}_{m,c}^{\rm inc,geom},
 \qquad
 \operatorname{Hol}_{m,c}^{\rm inc}:=\operatorname{Hol}_{m,c,g}^{\rm inc,spec}.
\end{equation}
Bajo $G^{V_m}$ se transforma por conjugación y bajo
$\prod_c\mathbb F_3^4$ permanece fija. Como $t_3$ tiene orden tres, el cálculo
funcional de esta variable de tres valores da la identidad literal
\begin{equation}\label{eq:pdf2-ym-owner-wc-curvature-functional}
 W_{c,g}:=\mathcal T_{m,g}^{-1}W_c^0
 =\mathbf1_{\{e\}}\!\left(\operatorname{Hol}_{m,c}^{\rm inc}\right)-\frac13.
\end{equation}
Se escribe $W_c:=W_{c,g}$ en el resto del capítulo.
Como $T_{m,b,g}$ actúa dentro del cociente del mismo collar, la holonomía espectral
es una función de su álgebra local de holonomías geométricas; no añade una variable
ni cambia el soporte. Las holonomías geométricas, usadas en $S_{m,g}$ y en el clover,
mantienen su distribución dependiente de $g$.
El proyector del miembro derecho es una función de clase local; por tanto $W_c$ no
es una coordenada auxiliar desacoplada, sino un observable de la holonomía de
incidencia de la misma restricción gauge.

Para un abierto acotado $O\subset\mathbb R^4$, sea
$\mathfrak A_{\rm YM}(O)$ el álgebra generada por funciones de clase acotadas de
\eqref{eq:pdf2-ym-owner-incidence-holonomy}, holonomías de plaqueta y lectores
clover con soporte en $O$. Se define el sector de curvatura física por
\begin{equation}\label{eq:pdf2-ym-owner-curvature-cyclic-sector}
 \mathscr H_{\rm YM}^{\rm curv}
 :=\overline{\bigcup_{O\Subset\mathbb R^4}
   \mathfrak A_{\rm YM}(O)\Omega_\infty}
 \subseteq\widehat{\mathscr H}_\infty^{\rm phys}.
\end{equation}
Por \eqref{eq:pdf2-ym-owner-wc-curvature-functional}, cada
$\widehat J_{m,\infty}W_c$ pertenece materialmente a este sector.
Las expectativas de coordenada envían cilindros de curvatura a cilindros de
curvatura; por tanto el semigrupo deja invariante
$\mathscr H_{\rm YM}^{\rm curv}$. Al ser autoadjunto, el sector es reductor. La
cota de \eqref{eq:pdf2-ym-owner-limit-gap} se restringe a él y el vector anterior
alcanza la igualdad:
\begin{equation}\label{eq:pdf2-ym-owner-curvature-sector-gap}
 D_{\rm curv,g}^{\rm YM}
 :=\left.D_{\infty,g}^{\rm YM}\right|_{\mathscr H_{\rm YM}^{\rm curv}},
 \qquad D_{\rm curv}^{\rm YM}:=D_{\rm curv,g}^{\rm YM},
 \qquad
 \inf\bigl(\operatorname{Spec}D_{\rm curv,g}^{\rm YM}\setminus\{0\}\bigr)
 =3\kappa_{\rm av}.
\end{equation}

\subsubsection{Sistema cofinal de representación y productos cruzados}

La inclusión $\widehat J_{\alpha\beta}$ conserva una coordenada ancestral; no debe
confundirse con el promedio de bloque que expresa un campo continuo. Para escribir
este segundo mapa sin alterar la genealogía, sea $\mathfrak P$ el conjunto dirigido
de pantallas celulares acotadas, cerrado bajo overlay, refinamiento nonádico y
transporte euclídeo. Para $\alpha\in\mathfrak P$ póngase
\begin{equation}\label{eq:pdf2-ym-owner-cell-step-space}
 V_\alpha:=\operatorname{span}\left\{
 e_{\alpha,c}:=|c|^{-1/2}\mathbf1_c:
 c\in\mathcal C_\alpha^{\rm cell}\right\}
 \subset L^2(\mathbb R^4),
 \qquad P_\alpha:L^2(\mathbb R^4)\to V_\alpha.
\end{equation}
Los overlays hacen que $P_\alpha h\to h$ para todo $h\in L^2$ a lo largo de
$\mathfrak P$. En el lado gauge se normaliza el testigo por
\begin{equation}\label{eq:pdf2-ym-owner-normalized-cell-witness}
 Z_{\alpha,c}:=\sqrt{\frac92}\,W_{\alpha,c},
 \qquad
 \langle Z_{\alpha,c},Z_{\alpha,d}\rangle=\delta_{cd},
 \qquad
 D_\alpha^{\rm YM}Z_{\alpha,c}=3\kappa_{\rm av}Z_{\alpha,c}.
\end{equation}
La segunda igualdad usa exactamente los momentos de
\eqref{eq:pdf2-ym-owner-wc-moments} y la factorización de coordenadas distintas.
Sea $\mathcal K_\alpha^{\rm cell}$ el espacio generado por esos vectores. Si
$\beta\succeq\alpha$, el entrelazador de representación, distinto de
$\widehat J_{\alpha\beta}$, se define generador a generador por
\begin{equation}\label{eq:pdf2-ym-owner-block-intertwiner}
 I_{\alpha\beta}^{\rm curv}Z_{\alpha,c}
 :=\sum_{\substack{d\in\mathcal C_\beta^{\rm cell}\\d\subset c}}
       \sqrt{\frac{|d|}{|c|}}\,Z_{\beta,d}.
\end{equation}
Como las subceldas son disjuntas y suman $|c|$, se obtiene, sin tomar límites,
\begin{equation}\label{eq:pdf2-ym-owner-block-intertwiner-identities}
 (I_{\alpha\beta}^{\rm curv})^*I_{\alpha\beta}^{\rm curv}=I,
 \quad
 I_{\beta\gamma}^{\rm curv}I_{\alpha\beta}^{\rm curv}
 =I_{\alpha\gamma}^{\rm curv},
 \quad
 D_\beta^{\rm YM}I_{\alpha\beta}^{\rm curv}
 =I_{\alpha\beta}^{\rm curv}D_\alpha^{\rm YM}.
\end{equation}
El relabelado celular conserva volúmenes y descendientes; para $a\in E(4)$ se tiene además
\begin{equation}\label{eq:pdf2-ym-owner-block-e4-naturality}
 U_\beta(a)I_{\alpha\beta}^{\rm curv}
 =I_{a\alpha,a\beta}^{\rm curv}U_\alpha(a).
\end{equation}

Para $h\in C_c^\infty(\mathbb R^4)$ defínase ahora el lector de bloque
\begin{equation}\label{eq:pdf2-ym-owner-smeared-witness}
 \Phi_\alpha(h):=
 \sum_{c\in\mathcal C_\alpha^{\rm cell}}
 \langle h,e_{\alpha,c}\rangle_{L^2}Z_{\alpha,c}.
\end{equation}
En una malla uniforme esta fórmula es
$\sqrt{9/2}\,a_\alpha^2\sum_c(P_\alpha h)(x_c)W_{\alpha,c}$; el promedio celular,
no una evaluación puntual, fija el producto cruzado. En efecto, si $\gamma$ refina
dos pantallas $\alpha$ y $\beta$, las definiciones anteriores dan la identidad
calculada
\begin{equation}\label{eq:pdf2-ym-owner-smeared-cross-gram}
 \left\langle
  I_{\alpha\gamma}^{\rm curv}\Phi_\alpha(h),
  I_{\beta\gamma}^{\rm curv}\Phi_\beta(k)
 \right\rangle
 =\langle P_\alpha h,P_\beta k\rangle_{L^2}.
\end{equation}
Por polarización, con $k=h$,
\begin{equation}\label{eq:pdf2-ym-owner-smeared-cauchy-derived}
 \left\|
  I_{\alpha\gamma}^{\rm curv}\Phi_\alpha(h)
  -I_{\beta\gamma}^{\rm curv}\Phi_\beta(h)
 \right\|_2^2
 =\|P_\alpha h-P_\beta h\|_{L^2}^2\longrightarrow0.
\end{equation}
Así la propiedad de Cauchy es consecuencia de la convergencia fuerte de
proyecciones, no una identidad adicional postulada. El colímite de
\eqref{eq:pdf2-ym-owner-block-intertwiner} contiene, por tanto,
\begin{equation}\label{eq:pdf2-ym-owner-smeared-gram}
 \Phi(h):=\lim_\alpha I_{\alpha,\infty}^{\rm curv}\Phi_\alpha(h),
 \qquad
 \langle\Phi(h),\Phi(k)\rangle=\langle h,k\rangle_{L^2},
 \qquad
 D_{\rm pub}^{\rm YM}\Phi(h)=3\kappa_{\rm av}\Phi(h).
\end{equation}

Esta representación no añade grados de libertad. En cada paso, la descomposición
$V_\beta=V_\alpha\oplus(V_\beta\ominus V_\alpha)$ tiene una base de Haar
nonádica. El orden exhaustivo de incidencias nuevas de
\eqref{eq:pdf2-ym-owner-exhaustive-index} empareja esa base, etiqueta por etiqueta,
con coordenadas $Z_{\beta,d}$ del factor de innovación, conservando el soporte de la
función de Haar. Las unitarias triangulares $\widehat\Gamma_\alpha$ pegan esos
emparejamientos y producen una isometría basada
\begin{equation}\label{eq:pdf2-ym-owner-publication-into-physical}
 \mathcal U_{\rm curv}:
 \varinjlim(\mathcal K_\alpha^{\rm cell},I_{\alpha\beta}^{\rm curv})
 \longrightarrow\mathscr H_{\rm YM}^{\rm curv},
 \quad
 D_{\rm curv}^{\rm YM}\mathcal U_{\rm curv}
 =\mathcal U_{\rm curv}D_{\rm pub}^{\rm YM},
 \quad U(a)\mathcal U_{\rm curv}=\mathcal U_{\rm curv}U_{L^2}(a).
\end{equation}
En adelante $\Phi(h)$ denota también su imagen por $\mathcal U_{\rm curv}$.
La compatibilidad de \eqref{eq:pdf2-ym-gibbs-transport-global} define
$\mathcal T_{\infty,g}$ en el límite. La representación física al acoplamiento fijado
no es una segunda isometría escogida sólo sobre el primer caos, sino la restricción
\begin{equation}\label{eq:pdf2-ym-full-publication-restriction}
 \mathcal U_{{\rm curv},g}
 :=\mathcal T_{\infty,g}^{-1}\mathcal U_{{\rm curv},0},
 \qquad
 \widetilde{\mathcal U}_g
 :=\mathcal T_{\infty,g}^{-1}:
 L^\infty(\widehat X_\infty,\widehat\mu_\infty^0)
 \longrightarrow
 L^\infty(\widehat X_\infty,\widehat\mu_{\infty,g}^{\rm YM}).
\end{equation}
Sobre la $*$-álgebra polinómica generada por holonomías, plaquetas, clover y
$\Phi(h)$ se verifican, generador a generador y después por linealidad,
\begin{equation}\label{eq:pdf2-ym-full-publication-products}
 \widetilde{\mathcal U}_g(AB)
 =\widetilde{\mathcal U}_g(A)\widetilde{\mathcal U}_g(B),
 \quad
 \widetilde{\mathcal U}_g(A^*)=\widetilde{\mathcal U}_g(A)^*,
 \quad
 \omega_g(\widetilde{\mathcal U}_gA)=\omega_0(A).
\end{equation}
Las igualdades de superficie, clover y Bianchi son preservadas por
\eqref{eq:pdf2-ym-transport-surface-bianchi}; por densidad monótona, la extensión
vale en toda el álgebra local acotada. Así
\eqref{eq:pdf2-ym-owner-smeared-cross-gram} es la restricción de un isomorfismo
probabilístico y algebraico completo, no un sustituto de él.

\subsubsection{Núcleo denso, continuidad fuerte y red local}

Sea $M_{\Phi(h)}$ el operador de multiplicación afiliado al álgebra local y
\begin{equation}\label{eq:pdf2-ym-owner-curvature-weyl}
 \mathsf W(h,t):=\exp\!\bigl(itM_{\Phi(h)}\bigr),
 \qquad h\in C_c^\infty(\mathbb R^4),\quad t\in\mathbb R.
\end{equation}
Se completa $\mathfrak A_{\rm YM}(O)$ con estos unitarios para
$\operatorname{supp}h\subset O$ y con las funciones de clase acotadas de las
holonomías de plaqueta. Denótese por $\mathfrak A_{\rm YM}(O)^{\rm sm}$ la
$*$-álgebra así generada con funciones de test suaves. El núcleo
\begin{equation}\label{eq:pdf2-ym-owner-smooth-core}
 \mathscr D_{\rm sm}:=\operatorname{span}\left\{
 A_1\cdots A_n\Omega_\infty:
 A_j\in\mathfrak A_{\rm YM}(O_j)^{\rm sm},\quad
 O_j\Subset\mathbb R^4\right\}
\end{equation}
es denso por la definición cíclica de
\eqref{eq:pdf2-ym-owner-curvature-cyclic-sector}. Las ecuaciones
\eqref{eq:pdf2-ym-owner-block-e4-naturality} y
\eqref{eq:pdf2-ym-owner-publication-into-physical} dan, para $a\in E(4)$,
\begin{equation}\label{eq:pdf2-ym-owner-field-e4-covariance}
 U(a)\Phi(h)=\Phi(a\!\cdot h),
 \qquad
 \|(U(a)-I)\Phi(h)\|=\|a\!\cdot h-h\|_{L^2}.
\end{equation}
Además, $|e^{ix}-e^{iy}|\le|x-y|$ y
\eqref{eq:pdf2-ym-owner-smeared-gram} implican
\begin{equation}\label{eq:pdf2-ym-owner-weyl-continuity}
 \|\bigl(\mathsf W(a\!\cdot h,t)-\mathsf W(h,t)\bigr)\Omega_\infty\|
 \le |t|\,\|a\!\cdot h-h\|_{L^2}.
\end{equation}
La continuidad se comprueba también para los generadores de holonomía del core;
\eqref{eq:pdf2-ym-owner-weyl-continuity} por sí sola no los incluye. Para una ruta poligonal orientada $\gamma$,
una plaqueta $p$ y $f\in C^1(G)^{\rm class}$ se ponen
\begin{equation}\label{eq:pdf2-ym-owner-holonomy-core-generators}
 A_{\gamma,f}:=f(\operatorname{Hol}_\gamma),
 \qquad A_{p,f}:=f(\operatorname{Hol}_p).
\end{equation}
En un overlay de $\gamma$ y $a\gamma$, con $a\in E(4)$, la ley de puente
\eqref{eq:pdf2-ym-owner-heat-bridge} cancela los incrementos comunes. Para cada
incremento restante $z$ se usa
$|f(xz)-f(x)|\le\|\nabla f\|_\infty d_G(z,e)$ y la cota de segundo momento del
kernel de calor
\(
 \int_Gd_G(z,e)^2k_\tau(z)\,dz\le C_G\tau
\).
Para la ley reponderada, $0<e^{-p/(4g^2)}\le1$ y la continuidad de $p$ en el
factor compacto dan
\begin{equation}\label{eq:pdf2-ym-owner-gibbs-heat-moment}
 \frac{\int_Gd_G(z,e)^2e^{-p(z)/(4g^2)}k_\tau(z)\,dz}
      {\int_Ge^{-p(z)/(4g^2)}k_\tau(z)\,dz}
 \le C_{G,g}\tau,
\end{equation}
uniformemente en las pantallas de un compacto y para $0<\tau\le\tau_0$.
Sumando los tiempos de los incrementos de la diferencia simétrica se obtiene la
estimación calculada
\begin{equation}\label{eq:pdf2-ym-owner-holonomy-e4-estimate}
 \begin{aligned}
 \|U(a)A_{\gamma,f}\Omega_\infty-A_{\gamma,f}\Omega_\infty\|_2^2
 &\le C_{G,g}\|\nabla f\|_\infty^2
     \tau(\gamma\triangle a\gamma),\\
 \|U(a)A_{p,f}\Omega_\infty-A_{p,f}\Omega_\infty\|_2^2
 &\le C_{G,g}\|\nabla f\|_\infty^2
     \tau(\partial p\triangle\partial(ap)).
 \end{aligned}
\end{equation}
Aquí la diferencia simétrica no se mide como dos curvas disjuntas: los conectores
de superficie la llenan por la cinta orientada mínima
$\Sigma(\gamma,a\gamma)$ y
\begin{equation}\label{eq:pdf2-ym-owner-bridge-time-ribbon}
 \tau(\gamma\triangle a\gamma)
 :=\sum_{p\subset\Sigma(\gamma,a\gamma)}\tau_p,
 \qquad
 \tau(\gamma\triangle a\gamma)
 \le C_\gamma d_{E(4)}(a,e),
\end{equation}
con la misma definición para las fronteras de plaqueta. La desigualdad sale de
sumar los tiempos aditivos de \eqref{eq:pdf2-ym-owner-heat-bridge} sobre una cinta
de anchura $d_{E(4)}(a,e)$ y longitud acotada por la de $\gamma$.
La medida de tiempo de puente del miembro derecho tiende así a cero cuando $a\to e$;
esto es exactamente la continuidad de los conectores de ruta de la imagen gauge.
Cada lector clover acotado $\chi(F_{\rm cl})$, $\chi\in C_b^1$, es función de una
suma finita de plaquetas transportadas; la misma cota, multiplicada por
$\|\chi'\|_\infty$ y por el número de términos, prueba también su continuidad.

Para funciones de clase meramente acotadas se hace visible la regularización que
entra en el core. Si $d b$ es Haar izquierdo en $E(4)$ y
$\eta\in C_c^\infty(E(4))$, se define el integral de Bochner
\begin{equation}\label{eq:pdf2-ym-owner-orbit-smearing}
 A[\eta]:=\int_{E(4)}\eta(b)U(b)AU(b)^*\,db,
 \qquad A\in\{A_{\gamma,f},A_{p,f},\chi(F_{{\rm cl}})\},
 \quad \chi\in C_b^1.
\end{equation}
Con $(L_a\eta)(b):=\eta(a^{-1}b)$, el cambio de variable $b\mapsto ab$ da
\begin{equation}\label{eq:pdf2-ym-owner-orbit-smearing-continuity}
 U(a)A[\eta]U(a)^*=A[L_a\eta],
 \qquad
 \|(A[L_a\eta]-A[\eta])\Omega_\infty\|
 \le\|A\|\,\|L_a\eta-\eta\|_{L^1(E(4))}\longrightarrow0.
\end{equation}
Se entiende desde ahora que los generadores de holonomía, plaqueta y clover de
$\mathfrak A_{\rm YM}(O)^{\rm sm}$ son los de
\eqref{eq:pdf2-ym-owner-holonomy-core-generators} con $f\in C^1$, o sus
regularizaciones \eqref{eq:pdf2-ym-owner-orbit-smearing}, y que el soporte barrido
permanece compacto dentro de $O$. Las aproximaciones de identidad en $E(4)$,
junto con \eqref{eq:pdf2-ym-owner-holonomy-e4-estimate}, recuperan los generadores
no regularizados en $L^2$; por ello no se reduce el sector cíclico.

Un telescopio de productos que usa
\eqref{eq:pdf2-ym-owner-weyl-continuity} y
\eqref{eq:pdf2-ym-owner-orbit-smearing-continuity} prueba la continuidad de todos
los generadores de $\mathscr D_{\rm sm}$; por
densidad y unitariedad,
\begin{equation}\label{eq:pdf2-ym-owner-strong-e4}
 \lim_{a\to e}\|(U(a)-I)\Psi\|=0
 \qquad(\Psi\in\mathscr H_{\rm YM}^{\rm curv}).
\end{equation}
Los soportes disjuntos usan factores disjuntos y conmutan. Queda así la red
\begin{equation}\label{eq:pdf2-ym-owner-local-net}
\begin{aligned}
 O_1\subseteq O_2&\Rightarrow
 \mathfrak A_{\rm YM}(O_1)\subseteq\mathfrak A_{\rm YM}(O_2),\\
 [\mathfrak A_{\rm YM}(O_1),\mathfrak A_{\rm YM}(O_2)]&=0
 \quad (O_1\cap O_2=\varnothing),\\
 U(a)\mathfrak A_{\rm YM}(O)U(a)^*&=\mathfrak A_{\rm YM}(aO).
\end{aligned}
\end{equation}
La identidad de Gram para $h\ne0$ y
\eqref{eq:pdf2-ym-owner-wc-curvature-functional} prueban que la red no es escalar.

\subsubsection{Axiomas de Osterwalder--Schrader y reconstrucción continua}

Para una dirección $\nu$, sea $\mathfrak A_{+,\nu}$ la clausura de la unión de
las álgebras anteriores con soporte compacto en $\{x_\nu>0\}$. El estado
estático, la ley de pilas y el operador de transferencia no se definen por
separado. En una pantalla $\alpha$, sea
$\widehat\mu_{\alpha,g}^{\rm YM}$ la marginal de
\eqref{eq:pdf2-ym-gibbs-measure}, sea
$\nu_{\alpha,g}:=(\pi_\alpha^{\rm phys})_*
\widehat\mu_{\alpha,g}^{\rm YM}$ su ley en el cociente y sea
$K_{\alpha,g}(t)=e^{-tD_{\alpha,g}^{\rm YM}}$ su kernel.
Para tiempos ordenados $t_1\le\cdots\le t_n$, la marginal euclídea de la
misma medida es
\begin{equation}\label{eq:pdf2-ym-owner-euclidean-time-marginal}
 \mathbb P_{\alpha,g}^{E,(\nu;t_1,\ldots,t_n)}
 :=(\operatorname{ev}_{t_1}^{(\nu)},\ldots,
    \operatorname{ev}_{t_n}^{(\nu)})_*
    \widehat\mu_{\alpha,g}^{\rm YM}.
\end{equation}
La identidad de desintegración que forma parte de
\eqref{eq:pdf2-g-gau-transfer-euclidean-dato} la calcula, no la reemplaza por
una cadena auxiliar:
\begin{equation}\label{eq:pdf2-ym-owner-euclidean-markov-identification}
 d\mathbb P_{\alpha,g}^{E,(\nu;t_1,\ldots,t_n)}
 =\nu_{\alpha,g}(dy_1)
   \prod_{j=1}^{n-1}K_{\alpha,g}
      (t_{j+1}-t_j;y_j,dy_{j+1}).
\end{equation}
Si $M_A$ denota multiplicación por un lector acotado $A$, sus correlaciones
son por tanto
\begin{equation}\label{eq:pdf2-ym-owner-schwinger-transfer}
 \begin{aligned}
 S_{\alpha,n}^{(\nu,g)}(A_1,t_1;\ldots;A_n,t_n)
 &:=\int\prod_{j=1}^n(\tau_{t_je_\nu}A_j)(x)\,
       d\widehat\mu_{\alpha,g}^{\rm YM}(x)\\
 &=\langle\mathbf1,
 M_{A_1}K_{\alpha,g}(t_2-t_1)M_{A_2}\cdots
 K_{\alpha,g}(t_n-t_{n-1})M_{A_n}\mathbf1
 \rangle_{L^2(\nu_{\alpha,g})}.
 \end{aligned}
\end{equation}
La ley de la pila estacionaria en esos mismos tiempos es esa marginal
euclídea, escrita mediante su desintegración:
\begin{equation}\label{eq:pdf2-ym-owner-markov-stack-law}
 d\mathbb P_{\alpha,g}^{(\nu;t_1,\ldots,t_n)}
 :=d\mathbb P_{\alpha,g}^{E,(\nu;t_1,\ldots,t_n)}
 =\nu_{\alpha,g}(dy_1)
   \prod_{j=1}^{n-1}K_{\alpha,g}(t_{j+1}-t_j;y_j,dy_{j+1}).
\end{equation}
Integrar sucesivamente $y_n,y_{n-1},\ldots,y_1$ da la identidad exacta
\begin{equation}\label{eq:pdf2-ym-owner-stack-transfer-identity}
 S_{\alpha,n}^{(\nu,g)}(A_1,t_1;\ldots;A_n,t_n)
 =\int\prod_{j=1}^nA_j(y_j)\,
   d\mathbb P_{\alpha,g}^{(\nu;t_1,\ldots,t_n)}(y).
\end{equation}
Si todos los tiempos coinciden, los kernels son la identidad y
\begin{equation}\label{eq:pdf2-ym-owner-static-is-stack}
 S_{\alpha,n}^{(\nu,g)}(A_1,t;\ldots;A_n,t)
 =\int A_1\cdots A_n\,d\nu_{\alpha,g}.
\end{equation}
Así el estado Schwinger estático es la marginal de tiempo igual de la misma
medida euclídea, y su desintegración es la ley Markov; no son dos funcionales
que luego se identifican por nombre. La
compatibilidad de $\widehat J_{\alpha,g}$ con medida y kernel hace compatible
\eqref{eq:pdf2-ym-owner-stack-transfer-identity} con todo overlay. Su límite define
\begin{equation}\label{eq:pdf2-ym-owner-schwinger-functionals}
 S_n^{\rm YM}(A_1,t_1;\ldots;A_n,t_n)
 :=\lim_{\alpha}S_{\alpha,n}^{(\nu,g)}
 (A_{1,\alpha},t_1;\ldots;A_{n,\alpha},t_n).
\end{equation}

Las inserciones de $\Phi$ se obtienen derivando en $t=0$ los unitarios
\eqref{eq:pdf2-ym-owner-curvature-weyl}. Puesto que los $Z_{\alpha,c}$ transportados
son centrados, independientes y uniformemente acotados, el lema de Hoeffding da,
nivel a nivel,
\begin{equation}\label{eq:pdf2-ym-owner-subgaussian-bound}
 \left\langle\Omega_\alpha,
 e^{t\Phi_\alpha(h)}\Omega_\alpha\right\rangle
 \le \exp\!\left(\frac9{16}t^2\|P_\alpha h\|_{L^2}^2\right).
\end{equation}
La cota pasa al límite por
\eqref{eq:pdf2-ym-owner-smeared-cauchy-derived}; todas las derivadas son continuas en
la topología de Schwartz y definen distribuciones temperadas.

La forma de ocho caras se prolonga primero a toda pila finita. Si $F$ depende de
las $N$ capas positivas de una pantalla $\alpha$, la reflexión de
\eqref{eq:pdf2-ym-owner-markov-stack-law} y la condición en la capa central dan
\begin{equation}\label{eq:pdf2-ym-owner-finite-stack-os}
 \begin{aligned}
 &\int \overline{F(\Theta_\nu y_{-N},\ldots,
                         \Theta_\nu y_{-1})}
          F(y_1,\ldots,y_N)\,d\Omega_{\alpha,N}^{(\nu,g)}(y)\\
 &\qquad=
 \int\nu_{\alpha,g}(dz)\left|
  \int K_{\alpha,g}(a_\alpha/2;z,dy_1)
       \prod_{j=1}^{N-1}K_{\alpha,g}(a_\alpha;y_j,dy_{j+1})
       F(y_1,\ldots,y_N)
 \right|^2\ge0.
 \end{aligned}
\end{equation}
Todo cilindro de semiespacio pertenece a una pila finita; por ello
\eqref{eq:pdf2-ym-owner-finite-stack-os} materializa la forma OS antes de la
clausura. En el cociente OS, si $F$ tiene inserciones en tiempos positivos, sea
$(\tau_sF)(A_1,t_1;\ldots;A_n,t_n):=
F(A_1,t_1+s;\ldots;A_n,t_n+s)$. Se define el operador de traslado por
\begin{equation}\label{eq:pdf2-ym-owner-os-transfer-operator}
 T_{\rm OS}^{(\nu)}(s)[F]:=[\tau_sF],
 \qquad s\ge0.
\end{equation}
Para cilindros $F,G$ de semiespacio, la igualdad euclídea--Markov
\eqref{eq:pdf2-ym-owner-euclidean-markov-identification} y la propiedad de
semigrupo dan generador a generador
\begin{equation}\label{eq:pdf2-ym-owner-os-markov-intertwining}
 \begin{aligned}
 \langle[F],T_{\rm OS}^{(\nu)}(s)[G]\rangle_{\rm OS}
 &=S_g^{\rm YM}((\Theta_\nu F)^*\tau_sG)\\
 &=\left\langle
    \mathcal V_{\infty,\nu}[F],
    e^{-sD_{\rm curv,g}^{\rm YM}}
    \mathcal V_{\infty,\nu}[G]
   \right\rangle .
 \end{aligned}
\end{equation}
La identidad muestra a la vez que las clases nulas permanecen nulas y que,
por densidad de los cilindros,
\begin{equation}\label{eq:pdf2-ym-owner-os-generator-identification}
 \mathcal V_{\infty,\nu}T_{\rm OS}^{(\nu)}(s)
 \mathcal V_{\infty,\nu}^{-1}
 =e^{-sD_{\rm curv,g}^{\rm YM}}.
\end{equation}
Ésta es la flecha explícita medida euclídea $\to$ desintegración $\to$
transferencia OS $\to$ Hamiltoniano.

\begin{proposition}[Paquete Osterwalder--Schrader continuo]
\label{prop:pdf2-ym-owner-continuum-os-package}
Los funcionales \eqref{eq:pdf2-ym-owner-schwinger-functionals} satisfacen:
\begin{enumerate}
 \item normalización, hermiticidad y simetría bosónica;
 \item covariancia euclídea y regularidad temperada;
 \item positividad por reflexión en cada una de las cuatro direcciones,
 \begin{equation}\label{eq:pdf2-ym-owner-continuum-reflection-positive}
  \sum_{i,j}\overline{a_i}a_j\,
  S^{\rm YM}\!\left((\Theta_\nu F_i)^*F_j\right)\ge0,
  \qquad F_i\in\mathfrak A_{+,\nu};
 \end{equation}
 \item localidad y cluster exponencial: si $A,B$ son locales centrados y una
 traslación de longitud $r$ separa sus soportes, entonces
 \begin{equation}\label{eq:pdf2-ym-owner-continuum-cluster}
  |S^{\rm YM}(A\,\tau_rB)|
  \le e^{-3\kappa_{\rm av}r}\,
       \|A\Omega_\infty\|\,\|B\Omega_\infty\|.
 \end{equation}
\end{enumerate}
La reconstrucción OS de cualquiera de las cuatro semirredes produce el mismo
espacio cíclico, el mismo vacío y el Hamiltoniano
$D_{\rm curv,g}^{\rm YM}$. La continuación OS produce campos de Wightman locales y una
representación de Poincaré cuyo espectro satisface
\begin{equation}\label{eq:pdf2-ym-owner-wightman-spectrum}
 \operatorname{Spec}(P^0,\mathbf P)\subseteq\overline V_+,
 \qquad P^0=D_{\rm curv,g}^{\rm YM}\ge0.
\end{equation}
\end{proposition}

\begin{proof}
Normalización y hermiticidad proceden del estado; simetría y localidad, de la
conmutatividad euclídea a soportes disjuntos. La covariancia y continuidad son
\eqref{eq:pdf2-ym-owner-field-e4-covariance}--
\eqref{eq:pdf2-ym-owner-strong-e4}; la regularidad es
\eqref{eq:pdf2-ym-owner-subgaussian-bound}. Para
$F_i$ cilíndricos, todos viven en un overlay y una pila finitos;
\eqref{eq:pdf2-ym-owner-finite-stack-os} y la polarización prueban
\eqref{eq:pdf2-ym-owner-continuum-reflection-positive}. La aproximación en norma OS
y Cauchy--Schwarz la extienden a $\mathfrak A_{+,\nu}$.

Tras una rotación euclídea, la separación espacial se convierte en tiempo positivo.
La identidad pila--transferencia
\eqref{eq:pdf2-ym-owner-stack-transfer-identity} y el operador
\eqref{eq:pdf2-ym-owner-os-transfer-operator} escriben la correlación centrada como
\[
 \langle A^*\Omega_\infty,
 e^{-rD_{\rm curv,g}^{\rm YM}}B\Omega_\infty\rangle.
\]
La cota espectral \eqref{eq:pdf2-ym-owner-curvature-sector-gap} da
\eqref{eq:pdf2-ym-owner-continuum-cluster}. Finalmente, los conectores
\eqref{eq:pdf2-ym-owner-os-connectors} forman un sistema dirigido; su unión es densa
por \eqref{eq:pdf2-ym-owner-smooth-core}. La identidad
\eqref{eq:pdf2-ym-owner-common-os-hamiltonian} identifica las cuatro clausuras y sus
generadores con $D_{\rm curv,g}^{\rm YM}$. Las hipótesis ya verificadas de regularidad,
covariancia, reflexión y cluster son las entradas de la continuación OS; su teorema de
reconstrucción da localidad de Wightman y
\eqref{eq:pdf2-ym-owner-wightman-spectrum}.
\end{proof}

\subsubsection{Reconstrucción cuatridimensional de la teoría de Yang--Mills}

La identificación no se apoya sólo en evaluar un clover sobre una conexión dada.
Sea $\mathcal R_G^{\rm loc}$ la $*$-álgebra de lectores formada por todas las
funciones de clase de las holonomías de incidencia y plaqueta, sus transportes a un
vértice base y las polarizaciones del lector $\mathcal P^{F^2}$. Sea
\begin{equation}\label{eq:pdf2-ym-owner-curvature-fibre}
 \mathcal E_G:=\Lambda^2(\mathbb R^4)^*\otimes\mathfrak g
\end{equation}
con el producto invariante. En cada celda, la polarización de
$\mathcal P^{F^2}$ da un Gram positivo $G_{\alpha,c}^{F}$ sobre
$\mathcal E_G$. Se divide por su kernel y se aplica su raíz positiva; el lector
normalizado resultante se denota $Y_{\alpha,c}(v)$ y verifica
\begin{equation}\label{eq:pdf2-ym-owner-curvature-cell-gram}
 \langle Y_{\alpha,c}(v),Y_{\alpha,d}(w)\rangle
 =\delta_{cd}\langle v,w\rangle_{\mathcal E_G}.
\end{equation}
Si $d\subset c$, sea $T_{d\leftarrow c}$ el transporte de marco ya presente en la
holonomía superficial. El entrelazador de curvatura completa es
\begin{equation}\label{eq:pdf2-ym-owner-curvature-intertwiner}
 I_{\alpha\beta}^{F}Y_{\alpha,c}(v)
 :=\sum_{d\subset c}\sqrt{\frac{|d|}{|c|}}\,
 Y_{\beta,d}\!\left(\operatorname{Ad}_{T_{d\leftarrow c}}v\right).
\end{equation}
La invariancia del producto de $\mathfrak g$, la disjunción de subceldas y la
composición de transportes prueban
\begin{equation}\label{eq:pdf2-ym-owner-curvature-intertwiner-identities}
 (I_{\alpha\beta}^{F})^*I_{\alpha\beta}^{F}=I,
 \qquad
 I_{\beta\gamma}^{F}I_{\alpha\beta}^{F}=I_{\alpha\gamma}^{F}.
\end{equation}
Para $h\in C_c^\infty(\mathbb R^4;\mathcal E_G)$ se pone
\begin{equation}\label{eq:pdf2-ym-owner-f-alpha}
 \mathbf F_\alpha(h):=
 \sum_{c\in\mathcal C_\alpha^{\rm cell}}
 Y_{\alpha,c}\!\left(
   |c|^{-1/2}\int_c h(x)\,d^4x\right).
\end{equation}
En un overlay común, el mismo cálculo de
\eqref{eq:pdf2-ym-owner-smeared-cross-gram} da ahora
\begin{equation}\label{eq:pdf2-ym-owner-curvature-cross-gram}
 \left\langle I_{\alpha\gamma}^{F}\mathbf F_\alpha(h),
 I_{\beta\gamma}^{F}\mathbf F_\beta(k)\right\rangle
 =\langle P_\alpha h,P_\beta k\rangle_{L^2(\mathcal E_G)}.
\end{equation}
Por tanto $\mathbf F_\alpha(h)$ es Cauchy por convergencia fuerte de las proyecciones,
y no por una hipótesis nueva. El mismo emparejamiento de Haar e innovaciones de
\eqref{eq:pdf2-ym-owner-publication-into-physical}, ahora aplicado a la base
ortonormalizada de $G_{\alpha,c}^{F}$, realiza el colímite mediante una isometría
local y gauge-covariante en
$L^2(\widehat X_\infty,\widehat\mu_{\infty,g}^{\rm YM})$. Su límite define la
distribución aleatoria
gauge-covariante
\begin{equation}\label{eq:pdf2-ym-owner-distributional-curvature}
\begin{aligned}
 \mathbf F_{\mu\nu}&\in
 \mathcal S'(\mathbb R^4;\mathfrak g)\ \widehat\otimes\
 L^2(\widehat X_\infty,\widehat\mu_{\infty,g}^{\rm YM}),\\
 U(a)\mathbf F(h)U(a)^*&=\mathbf F(a\!\cdot h),\\
 u\mathbf F(h)u^{-1}&=\mathbf F(\operatorname{Ad}_u h),
 \qquad a\in E(4),\ u\in G.
\end{aligned}
\end{equation}
Su canal de clase espectral de orden tres es precisamente $\Phi$. La relación
superficial finita
\begin{equation}\label{eq:pdf2-ym-owner-surface-product}
 \operatorname{Hol}^{(\alpha)}_p
 =\prod_{p'\subset p}^{\longrightarrow}
  T_{p'}\operatorname{Hol}^{(\beta)}_{p'}T_{p'}^{-1}
\end{equation}
pasa a la clausura cofinal porque todos sus productos cruzados se calculan en un
overlay. En particular, el lector clover aleatorio, y no sólo su evaluación en una
conexión suave, satisface la identidad de pequeños lazos
\begin{equation}\label{eq:pdf2-ym-owner-holonomy-curvature-relation}
 \mathbf F(h)=L^2\!\mathord{-}\!\lim_{\alpha}
 a_\alpha^4\sum_x
 \left\langle(P_\alpha h)(x),F_{{\rm cl},\alpha}(x)\right\rangle,
 \qquad
 \delta_{\rm gauge}S_n^{\rm YM}=0,
 \qquad
 \operatorname{Alt}_{\lambda\mu\nu}D_\lambda\mathbf F_{\mu\nu}=0.
\end{equation}
La primera igualdad es \eqref{eq:pdf2-ym-owner-curvature-cross-gram} escrita en el
representante clover; la segunda es la identidad de Ward
\eqref{eq:pdf2-ym-ward-finite} pasada al límite, y la
tercera es la cancelación de la frontera de una frontera en el producto superficial.
Ninguna usa el término de masa.

Para evitar definir el blanco por la conclusión que se quiere obtener, sea
$\mathbf{YM}_4(G,g)$ la clase de quíntuplas locales
$(\mathfrak A,\mu_g,K_g,\operatorname{Hol},\mathbf F)$ que verifican las siguientes
ecuaciones comprobables: (i) la densidad de Gibbs es
\eqref{eq:pdf2-ym-gibbs-measure}; (ii) $K_g$ es Markov, reversible y satisface
\eqref{eq:pdf2-ym-schwinger-dyson-finite}--\eqref{eq:pdf2-ym-ward-finite};
(iii) holonomía, clover y curvatura satisfacen
\eqref{eq:pdf2-ym-transport-surface-bianchi} y
\eqref{eq:pdf2-ym-owner-holonomy-curvature-relation}; y (iv) sus funcionales son la
misma ley de pilas de \eqref{eq:pdf2-ym-owner-stack-transfer-identity} y cumplen el
paquete OS\@. El objeto construido es, componente a componente,
\begin{equation}\label{eq:pdf2-ym-owner-qym-object}
\begin{aligned}
 \mathfrak Q_{\rm YM}^{(4)}(G,g):={}&
 \bigl(\{\mathfrak A_{\rm YM}(O)\}_O,
 S_g^{\rm YM},U,\Theta,\Omega_{\infty,g},\\
 &\widehat\mu_{\infty,g}^{\rm YM},
 \{S_{m,g},\mathscr S_{m,g}\}_m,K_g,
 \operatorname{Hol},\mathbf F,\\
 &\mathscr H_{\rm OS},D_{\rm curv,g}^{\rm YM}\bigr)
 \in\mathbf{YM}_4(G,g).
\end{aligned}
\end{equation}
El mapa de representación no omite ninguna de esas componentes:
\begin{equation}\label{eq:pdf2-ym-owner-typed-publication-map}
 \begin{aligned}
 \mathcal P_{{\rm YM},g}:\mathfrak G_{\rm gau}(G)&\longrightarrow
 \mathbf{YM}_4(G,g),\\
 (\widehat\mu^0,\mathbb P^{\rm phys},D^0,\Theta,W,
   \operatorname{Hol},\mathcal P^{F^2})
 &\longmapsto
 (\widehat\mu_g^{\rm YM},S_g,\mathscr S_g,K_g,\mathbb P_g^{\rm phys},
  S_g^{\rm YM},\mathfrak A_{\rm YM},D_{\rm curv,g}^{\rm YM},
  \mathbf F,\operatorname{Hol}).
 \end{aligned}
\end{equation}
Los $S_n^{\rm YM}$ son así exactamente las funciones de Schwinger de la medida
dependiente de $g$, la ley Markov y el operador de transferencia ya construidos.

Queda por comprobar que $S_{m,g}$ es la acción de Yang--Mills y no sólo una
función con ese nombre. Sea
$A\in C_c^3(\mathbb R^4;T^*\mathbb R^4\otimes\mathfrak g)$ una conexión suave en
una coordenatización donde el logaritmo de los lazos pequeños es único. Para la plaqueta
$p_{\mu\nu}(x)$, la expansión ordenada de Magnus da
\begin{equation}\label{eq:pdf2-ym-owner-plaquette-magnus}
 \log\operatorname{Hol}_{p_{\mu\nu}(x)}(A)
 =a_m^2F_{A,\mu\nu}(x)
  +\frac{a_m^3}{2}(D_\mu+D_\nu)F_{A,\mu\nu}(x)
  +R_{m,\mu\nu}(x),
\end{equation}
con
$\|R_{m,\mu\nu}(x)\|\le C_G a_m^4
(\|D^2F_A\|_{\infty,p}+\|F_A\|_{\infty,p}^2)$.
El promedio de las cuatro orientaciones cancela el término impar y define
\begin{equation}\label{eq:pdf2-ym-owner-clover}
 F_{{\rm cl},m}(x)
 :=\frac1{4a_m^2}\sum_{p\ni x}
  \operatorname{Ad}_{T_p}\log\operatorname{Hol}_p,
 \qquad
 \|F_{{\rm cl},m}(x)-F_A(x)\|
 \le C_G a_m^2(\|D^2F_A\|_\infty+\|F_A\|_\infty^2).
\end{equation}
La definición de la raíz
\eqref{eq:pdf2-ym-action-root} es precisamente la polarización de esos seis
componentes. Sustituyendo \eqref{eq:pdf2-ym-owner-plaquette-magnus} en cada
generador de esa polarización se calcula
\begin{equation}\label{eq:pdf2-ym-owner-action-root-smooth-chart}
 \xi_{m,b}(A)
 =a_m^2\operatorname{vec}_{\mu<\nu}F_{A,\mu\nu}(x_b)
  +r_{m,b}(A),
 \qquad
 \|r_{m,b}(A)\|\le C_{G,A}a_m^4.
\end{equation}
Los términos cruzados impares se cancelan por las cuatro orientaciones del clover;
los pares están incluidos una sola vez por
\eqref{eq:pdf2-ym-action-factor-partition}. Por suma de Riemann y
Cauchy--Schwarz,
\begin{equation}\label{eq:pdf2-ym-owner-action-error}
 \left|S_{m,g}(A)-\frac1{4g^2}
  a_m^4\sum_x\|F_{{\rm cl},m}(x)\|^2\right|
 \le C_{G,A,g}a_m^2,
\end{equation}
y, usando la cota de \eqref{eq:pdf2-ym-owner-clover},
\begin{equation}\label{eq:pdf2-ym-owner-f2}
 S_{m,g}(A)\longrightarrow
 S_{{\rm YM},g}(A):=\frac1{4g^2}
 \int_{\mathbb R^4}\|F_A(x)\|^2\,d^4x.
\end{equation}
Las ecuaciones de Schwinger--Dyson y Ward ya se calcularon sobre la medida que usa
esta acción. La brecha, en cambio, se deduce después por la conjugación unitaria de
\eqref{eq:pdf2-ym-g-dynamics} y el testigo
\eqref{eq:pdf2-ym-g-gap-before-limit}; no fue usada para escoger $g$, $S_{m,g}$ ni
$\widehat\mu_{m,g}^{\rm YM}$. Éste es el enlace no circular entre acción, ley y gap.

\begin{theorem}[Cierre Yang--Mills para todo grupo de Lie compacto, conectado y simple]
\label{thm:pdf2-ym-cierre-global}
La estructura completa $\mathfrak G_{\rm gau}$ construye una colección proyectiva
local para todo grupo de Lie compacto, conectado y simple. El mapa
\eqref{eq:pdf2-ym-owner-typed-publication-map} expresa una teoría cuántica de
Yang--Mills no trivial sobre $\mathbb R^4$: sus funcionales satisfacen el paquete OS
continuo y su continuación satisface los axiomas locales de Wightman por la
proposición~\ref{prop:pdf2-ym-owner-continuum-os-package}; sus
holonomías y su curvatura
distribucional satisfacen
\eqref{eq:pdf2-ym-owner-holonomy-curvature-relation}, y las cuatro reconstrucciones
coinducen el mismo Hamiltoniano $D_{\infty,g}^{\rm YM}$, con vacío simple. Su ley es
$\widehat\mu_{\infty,g}^{\rm YM}$, construida por la colección compatible de acciones
locales $\{S_{m,g},\mathscr S_{m,g}\}_m$, y
satisface las ecuaciones Schwinger--Dyson y Ward
\eqref{eq:pdf2-ym-schwinger-dyson-finite}--\eqref{eq:pdf2-ym-ward-finite}. El sector
cíclico de curvatura
\eqref{eq:pdf2-ym-owner-curvature-cyclic-sector} es reductor, contiene el testigo
\eqref{eq:pdf2-ym-owner-wc-curvature-functional} y posee
\begin{equation}\label{eq:pdf2-ym-final-gap}
 \boxed{\Delta_{\rm YM}^{\rm HMT}=3\kappa_{\rm av}>0},
 \qquad
 m_{\rm gap}^{\rm HMT}
 =\frac{\hbar_{\rm int}}{c_{\rm int}}\Delta_{\rm YM}^{\rm HMT}>0.
\end{equation}
La red local, la acción fuerte de $E(4)$, las distribuciones de Schwinger, la
holonomía y $\mathbf F$ proceden del mismo estado completo. El lector exacto
$\mathcal P^{F^2}$ define la densidad de la ley; el clover es su coordenatización suave y
\eqref{eq:pdf2-ym-owner-f2} identifica su acción clásica. La brecha pertenece al sector colectivo
gauge-invariante de curvatura y no asigna una masa elemental al gluón.
\end{theorem}

\begin{proof}
La proyectividad cinemática es \eqref{eq:pdf2-ym-owner-projectivity}; la raíz
\eqref{eq:pdf2-ym-action-root}, el cociclo
\eqref{eq:pdf2-ym-action-cocycle} y la densidad
\eqref{eq:pdf2-ym-gibbs-measure} construyen primero la acción y la medida a
acoplamiento $g$. El transporte factor a factor
\eqref{eq:pdf2-ym-gibbs-transport-factor}--
\eqref{eq:pdf2-ym-full-algebra-identities} prueba proyectividad, preservación del
estado, productos y $*$; 
\eqref{eq:pdf2-ym-transport-surface-bianchi} descarga clover y Bianchi. El índice exhaustivo
\eqref{eq:pdf2-ym-owner-exhaustive-index}, la forma
\eqref{eq:pdf2-ym-owner-full-form} y el circuito
\eqref{eq:pdf2-ym-owner-local-circuit} construyen el generador común. La reducción
física y su kernel cociente son
\eqref{eq:pdf2-ym-owner-physical-reduction}--
\eqref{eq:pdf2-ym-owner-D}. La conjugación
\eqref{eq:pdf2-ym-g-dynamics} los transporta a la medida dependiente de $g$ sin
alterar su espectro. La descomposición de Hoeffding prueba el vacío y la cota, y
\eqref{eq:pdf2-ym-owner-wc-invariance}--
\eqref{eq:pdf2-ym-owner-wc-eigenvalue} exhiben un vector físico que alcanza el
umbral. El entrelazamiento nonádico transporta forma, vacío y testigo al límite, donde
\eqref{eq:pdf2-ym-owner-limit-gap}--
\eqref{eq:pdf2-ym-owner-exact-gap} fijan la brecha exacta. Las cuatro
factorizaciones \eqref{eq:pdf2-ym-owner-four-os}, la inclusión de rangos
\eqref{eq:pdf2-ym-owner-os-range-inclusion} y los conectores
\eqref{eq:pdf2-ym-owner-os-connectors} identifican sus reconstrucciones con ese
mismo generador. La holonomía
\eqref{eq:pdf2-ym-owner-incidence-holonomy} coloca materialmente $W_c$ en el sector
de curvatura, donde \eqref{eq:pdf2-ym-owner-curvature-sector-gap} prueba la brecha
exacta. El sistema \eqref{eq:pdf2-ym-owner-block-intertwiner} controla los productos
cruzados mediante \eqref{eq:pdf2-ym-owner-smeared-cross-gram};
\eqref{eq:pdf2-ym-full-publication-restriction}--
\eqref{eq:pdf2-ym-full-publication-products} elevan ese primer caos al isomorfismo
de toda el álgebra. La convergencia fuerte de $P_\alpha$ produce $\Phi$; las cotas
\eqref{eq:pdf2-ym-owner-holonomy-e4-estimate} y
\eqref{eq:pdf2-ym-owner-orbit-smearing-continuity} prueban además la continuidad
$E(4)$ de holonomías, plaquetas y clover. La desintegración euclídea
\eqref{eq:pdf2-ym-owner-euclidean-markov-identification}, la identidad exacta
\eqref{eq:pdf2-ym-owner-stack-transfer-identity} y el entrelazamiento
\eqref{eq:pdf2-ym-owner-os-markov-intertwining} enlazan la misma medida,
estado Schwinger, pila Markov, transferencia OS y generador del gap. La
proposición~\ref{prop:pdf2-ym-owner-continuum-os-package} prolonga las
cuatro formas finitas a las
semirredes completas, prueba regularidad y cluster y reconstruye el Hamiltoniano. Por
último, \eqref{eq:pdf2-ym-owner-distributional-curvature}--
\eqref{eq:pdf2-ym-owner-typed-publication-map} identifican esos mismos funcionales
como los Schwinger de la teoría gauge local; las integraciones por partes
\eqref{eq:pdf2-ym-schwinger-dyson-finite}--\eqref{eq:pdf2-ym-ward-finite} prueban
Schwinger--Dyson y Ward. Finalmente,
\eqref{eq:pdf2-ym-owner-plaquette-magnus}--\eqref{eq:pdf2-ym-owner-f2} identifican
la acción clásica antes de aplicar la equivalencia espectral. La cadena no usa el
gap para escoger la medida.
\end{proof}

\begin{falsifier}[Colapso de la dinámica producto]
Si se reemplaza $D_m^{\rm prod}$ por $D_m^{\rm fib}$, todos los sectores no
constantes reciben el mismo autovalor. Un vector de Hoeffding con $|S|=2$ debería
tener energía $6\kappa_{\rm av}$ por \eqref{eq:pdf2-ym-hoeffding}, pero el
semigrupo simple le asignaría $3\kappa_{\rm av}$. La identificación de ambos
generadores es falsa.
\end{falsifier}

\begin{falsifier}[Confusión entre inclusión genealógica y promedio de bloque]
Si se sustituye $I_{\alpha\beta}^{\rm curv}$ por
$\widehat J_{\alpha\beta}$ en
\eqref{eq:pdf2-ym-owner-smeared-cross-gram}, una coordenada gruesa persiste con
coeficiente uno mientras el coeficiente fino cambia por el cociente de escalas. Para
$a_{m+1}=a_m/9$, el producto cruzado ancestral lleva un factor $1/81$ y no converge
a la norma diagonal. La identidad
\eqref{eq:pdf2-ym-owner-smeared-cauchy-derived} falla; por eso los dos mapas se
mantienen separados.
\end{falsifier}

\begin{falsifier}[Ablación de la fibra de incidencia]
Si se sustituye $\widehat X_m$ por su marginal de enlaces, se borra $q_c$ y con ella
el observable físico $W_c$. El mapa de olvido no refleja la subálgebra gauge completa;
por tanto no puede utilizarse para negar el autovalor del objeto fuente. Objeto
sustituido; conclusión no transferible.
\end{falsifier}

\begin{falsifier}[Sustitución del lector exacto por su coordenatización suave]
Si se usa directamente el término asintótico de
\eqref{eq:pdf2-ym-owner-f2} como densidad en una pantalla finita, se borra el resto
controlado de \eqref{eq:pdf2-ym-owner-action-error} y se pierde el cociclo exacto
\eqref{eq:pdf2-ym-action-cocycle}. La densidad finita gobernante es
$Z_{m,g}^{-1}e^{-\mathcal P_m^{F^2}/(4g^2)}$; el clover y el límite suave la
identifican con la acción clásica, pero no la sustituyen.
\end{falsifier}

\paragraph{Articulación de las demostraciones.}
La separación entre fibra simple y generador producto, la medida común, la compresión
gauge, las cuatro condiciones de Osterwalder--Schrader, la forma límite y el testigo físico son
resultados exactos previamente establecidos en el dominio gauge
regular HMT. La desintegración con tiempos no uniformes, el índice
exhaustivo de expectativas, el kernel cociente, la matriz $B^+$, la acción explícita
en el colímite, el sistema cofinal de bloque, el paquete OS continuo, la curvatura
distribucional y la holonomía de incidencia son
formalizaciones explícitas: determinan los dominios y demuestran los enlaces antes
enunciados, manteniendo la brecha como resultado. Su composición completa
establece el teorema precedente.
La raíz positiva del lector $F^2$, su densidad Gibbs, el transporte local compatible,
las ecuaciones Schwinger--Dyson/Ward, el isomorfismo algebraico completo, las cotas
de continuidad de holonomías y la identidad pila--transferencia son
formalizaciones explícitas de las cuatro correspondencias. La identificación
entre el observable clover y la acción clásica se obtiene como consecuencia
de la expansión de Magnus.
